\documentclass{article}

% ============================================================
%  NeurIPS 2026 — submission (anonymous, with line numbers)
%  Switch to [main, final] after acceptance for camera-ready.
% ============================================================
\PassOptionsToPackage{numbers, compress}{natbib}
\usepackage{neurips_2026}          % default = main track, anonymous
% \usepackage[main, final]{neurips_2026}  % camera-ready

% ---- standard packages (from NeurIPS template) ----
\usepackage[utf8]{inputenc}
\usepackage[T1]{fontenc}
\usepackage{hyperref}
\usepackage{url}
\usepackage{booktabs}
\usepackage{amsfonts}
\usepackage{amsmath}
\usepackage{amssymb}
\usepackage{nicefrac}
\usepackage{microtype}
\usepackage{graphicx}

% ---- additional packages needed by our sections ----
\usepackage{tabularx}       % for Table 1 (simulator comparison)
\usepackage{enumitem}        % better enumerate control
\usepackage{soul}            % \hl{} highlighting (used in TODOs)

% ---- checkmark command for the comparison table ----
\usepackage{pifont}          % provides \ding{51} (✓) and \ding{55} (✗)
\newcommand{\cmark}{\ding{51}}

% ---- review / advisor annotation commands (suppress for camera-ready) ----
% ============================================================
% Review / Advisor Comment Commands for SceneFactory manuscript
% Usage:  % ============================================================
% Review / Advisor Comment Commands for SceneFactory manuscript
% Usage:  % ============================================================
% Review / Advisor Comment Commands for SceneFactory manuscript
% Usage:  \input{reviewcommands}  in preamble of main.tex
%
% NOTE: neurips_2026.sty already loads xcolor, so we do NOT
%       reload it here.  We also skip loading marginnote since
%       NeurIPS geometry is tight.
% ============================================================

% ---- extra packages for review annotations ----
\usepackage{tcolorbox}
\tcbuselibrary{skins,breakable}

% ---- colour palette -----------------------------------------------------
\definecolor{chatBlue}{HTML}{1565C0}
\definecolor{chatGreen}{HTML}{2E7D32}
\definecolor{chatOrange}{HTML}{E65100}
\definecolor{chatRed}{HTML}{C62828}
\definecolor{chatPurple}{HTML}{6A1B9A}
\definecolor{chatTeal}{HTML}{00838F}
\newcommand{\yicheng}[1]{\textcolor{brown}{#1}}

% ---- inline commands ----------------------------------------------------

\newcommand{\chatcomment}[1]{%
  {\color{chatBlue}\sffamily\bfseries [Comment:\ #1]}%
}
\newcommand{\chatrewrite}[1]{%
  {\color{chatGreen}\sffamily\bfseries [Rewrite:\ #1]}%
}
\newcommand{\chatcite}[1]{%
  {\color{chatOrange}\sffamily\bfseries [Cite:\ #1]}%
}
\newcommand{\chatfactcheck}[1]{%
  {\color{chatRed}\sffamily\bfseries [Fact-check:\ #1]}%
}
\newcommand{\chatstructure}[1]{%
  {\color{chatPurple}\sffamily\bfseries [Structure:\ #1]}%
}
\newcommand{\chattodo}[1]{%
  {\color{chatTeal}\sffamily\bfseries [TODO:\ #1]}%
}

% ---- block-level box for longer advisor notes ---------------------------
\newtcolorbox{chatbox}[1][Advisor Note]{%
  colback=chatBlue!5!white,
  colframe=chatBlue!70!black,
  fonttitle=\sffamily\bfseries,
  title={#1},
  breakable,
  left=4pt, right=4pt, top=4pt, bottom=4pt,
}

% ---- toggle: uncomment ALL lines below to hide review markup -----------
% \renewcommand{\chatcomment}[1]{}
% \renewcommand{\chatrewrite}[1]{}
% \renewcommand{\chatcite}[1]{}
% \renewcommand{\chatfactcheck}[1]{}
% \renewcommand{\chatstructure}[1]{}
% \renewcommand{\chattodo}[1]{}
% \renewcommand{\yicheng}[1]{#1}
% \newtcolorbox{chatbox}[1][]{blank, breakable, before skip=0pt, after skip=0pt}
  in preamble of main.tex
%
% NOTE: neurips_2026.sty already loads xcolor, so we do NOT
%       reload it here.  We also skip loading marginnote since
%       NeurIPS geometry is tight.
% ============================================================

% ---- extra packages for review annotations ----
\usepackage{tcolorbox}
\tcbuselibrary{skins,breakable}

% ---- colour palette -----------------------------------------------------
\definecolor{chatBlue}{HTML}{1565C0}
\definecolor{chatGreen}{HTML}{2E7D32}
\definecolor{chatOrange}{HTML}{E65100}
\definecolor{chatRed}{HTML}{C62828}
\definecolor{chatPurple}{HTML}{6A1B9A}
\definecolor{chatTeal}{HTML}{00838F}
\newcommand{\yicheng}[1]{\textcolor{brown}{#1}}

% ---- inline commands ----------------------------------------------------

\newcommand{\chatcomment}[1]{%
  {\color{chatBlue}\sffamily\bfseries [Comment:\ #1]}%
}
\newcommand{\chatrewrite}[1]{%
  {\color{chatGreen}\sffamily\bfseries [Rewrite:\ #1]}%
}
\newcommand{\chatcite}[1]{%
  {\color{chatOrange}\sffamily\bfseries [Cite:\ #1]}%
}
\newcommand{\chatfactcheck}[1]{%
  {\color{chatRed}\sffamily\bfseries [Fact-check:\ #1]}%
}
\newcommand{\chatstructure}[1]{%
  {\color{chatPurple}\sffamily\bfseries [Structure:\ #1]}%
}
\newcommand{\chattodo}[1]{%
  {\color{chatTeal}\sffamily\bfseries [TODO:\ #1]}%
}

% ---- block-level box for longer advisor notes ---------------------------
\newtcolorbox{chatbox}[1][Advisor Note]{%
  colback=chatBlue!5!white,
  colframe=chatBlue!70!black,
  fonttitle=\sffamily\bfseries,
  title={#1},
  breakable,
  left=4pt, right=4pt, top=4pt, bottom=4pt,
}

% ---- toggle: uncomment ALL lines below to hide review markup -----------
% \renewcommand{\chatcomment}[1]{}
% \renewcommand{\chatrewrite}[1]{}
% \renewcommand{\chatcite}[1]{}
% \renewcommand{\chatfactcheck}[1]{}
% \renewcommand{\chatstructure}[1]{}
% \renewcommand{\chattodo}[1]{}
% \renewcommand{\yicheng}[1]{#1}
% \newtcolorbox{chatbox}[1][]{blank, breakable, before skip=0pt, after skip=0pt}
  in preamble of main.tex
%
% NOTE: neurips_2026.sty already loads xcolor, so we do NOT
%       reload it here.  We also skip loading marginnote since
%       NeurIPS geometry is tight.
% ============================================================

% ---- extra packages for review annotations ----
\usepackage{tcolorbox}
\tcbuselibrary{skins,breakable}

% ---- colour palette -----------------------------------------------------
\definecolor{chatBlue}{HTML}{1565C0}
\definecolor{chatGreen}{HTML}{2E7D32}
\definecolor{chatOrange}{HTML}{E65100}
\definecolor{chatRed}{HTML}{C62828}
\definecolor{chatPurple}{HTML}{6A1B9A}
\definecolor{chatTeal}{HTML}{00838F}
\newcommand{\yicheng}[1]{\textcolor{brown}{#1}}

% ---- inline commands ----------------------------------------------------

\newcommand{\chatcomment}[1]{%
  {\color{chatBlue}\sffamily\bfseries [Comment:\ #1]}%
}
\newcommand{\chatrewrite}[1]{%
  {\color{chatGreen}\sffamily\bfseries [Rewrite:\ #1]}%
}
\newcommand{\chatcite}[1]{%
  {\color{chatOrange}\sffamily\bfseries [Cite:\ #1]}%
}
\newcommand{\chatfactcheck}[1]{%
  {\color{chatRed}\sffamily\bfseries [Fact-check:\ #1]}%
}
\newcommand{\chatstructure}[1]{%
  {\color{chatPurple}\sffamily\bfseries [Structure:\ #1]}%
}
\newcommand{\chattodo}[1]{%
  {\color{chatTeal}\sffamily\bfseries [TODO:\ #1]}%
}

% ---- block-level box for longer advisor notes ---------------------------
\newtcolorbox{chatbox}[1][Advisor Note]{%
  colback=chatBlue!5!white,
  colframe=chatBlue!70!black,
  fonttitle=\sffamily\bfseries,
  title={#1},
  breakable,
  left=4pt, right=4pt, top=4pt, bottom=4pt,
}

% ---- toggle: uncomment ALL lines below to hide review markup -----------
% \renewcommand{\chatcomment}[1]{}
% \renewcommand{\chatrewrite}[1]{}
% \renewcommand{\chatcite}[1]{}
% \renewcommand{\chatfactcheck}[1]{}
% \renewcommand{\chatstructure}[1]{}
% \renewcommand{\chattodo}[1]{}
% \renewcommand{\yicheng}[1]{#1}
% \newtcolorbox{chatbox}[1][]{blank, breakable, before skip=0pt, after skip=0pt}


% ---- graphics path (add figures/ when you have images) ----
\graphicspath{{figures/}}

% ============================================================
\title{SceneFactory: GPU-Accelerated Multi-Agent Driving Simulation\\
with Physics-Based Vehicle Dynamics}

\author{%
  Anonymous Author(s)
  % ---- Uncomment for camera-ready: ----
  % Yicheng Zhu\thanks{Corresponding author.} \\
  % Kate Gleason College of Engineering\\
  % Rochester Institute of Technology\\
  % Rochester, NY 14623\\
  % \texttt{yz8733@g.rit.edu} \\
  % \And
  % Zilin Bian \\
  % Rochester Institute of Technology / University of Pittsburgh\\
  % Rochester, NY 14623 \\
  % \texttt{zxbite@rit.edu} \\
}

\begin{document}

\maketitle

% ============================================================
%  Abstract
% ============================================================
\begin{abstract}
Training and evaluating autonomous driving agents at scale requires simulators that combine high physics fidelity with high computational throughput.
Existing platforms force a trade-off: physically realistic simulators such as CARLA and MetaDrive model vehicles as non-vectorized PhysX or Bullet articulations with per-agent Python control loops, limiting them to one scene at a time; GPU-accelerated systems such as Waymax and GPUDrive achieve massive parallelism by abandoning rigid-body physics entirely in favor of bicycle or kinematic models.
We present \textbf{SceneFactory}, a procedural scene construction and multi-agent training platform built on NVIDIA Isaac Sim and Isaac Lab that closes this gap.
At its core is a \emph{fully vectorized articulated PhysX vehicle}: joint states, observations, rewards, and policy inference are computed entirely through batched GPU tensor operations, eliminating the per-agent loops that bottleneck existing physics-based simulators and achieving up to \textbf{127$\times$} higher throughput than a non-vectorized PhysX baseline on the same hardware.
SceneFactory ingests road topologies from the Waymo Open Motion Dataset, procedurally constructs multi-world USD environments, and orchestrates the parallel simulation of hundreds of these articulated vehicles across dozens of independent road worlds on a single GPU.
A weather-to-friction module maps precipitation and road surface type to effective friction coefficients applied directly to the physics engine.
We train navigation policies via PPO and report throughput and policy performance.
To our knowledge, SceneFactory is the first system that simultaneously provides GPU-parallel multi-world execution, multi-agent support, full rigid-body physics with vectorized articulated vehicles, and weather-aware friction modeling.
\end{abstract}

% ============================================================
%  Paper body
% ============================================================
% ============================================================
%  sections/introduction.tex  —  Introduction (merged & cleaned)
%  Sources: SOC draft + sec/1_intro.tex (previous writeup)
% ============================================================
\section{Introduction}
\label{sec:introduction}
%Autonomous driving is something that people are actively exploring, and there existws a lot of simulators for autonomous vehicles. Why researchers need simulators for AV? Because we surely need a system that is not real world, to try the trained jdriving algorithms, as the real world deployment is costly, it's impossible to try things in scale in the real world. Local transportation departments are going to make a lot of administrative decisions every day, especially for urban areas. I should take some data about how much traffic decisions are made in a day, and there must be researches about the disruptive effect to traffic from those decisions. Traffic engineers spent decades to make traffic simulators that work on both microscopic level and macroscopic level. ON one hand, we have SUMO, which is an opensource traffic simulator, that is widely used in digital twin technology. SUMO is a quite macroscope simjulator, which simulates how the traffic system will respond to a certain change to the system, such as road closure, lane closure. However, there are plenty of things that are simplified in SUMO, such as vehicle dynamics. Differetn weathers will have effects on how cars dynamics change, and how the drivers behaviors change. Those changes cannot be captured by SUMO's modeling of driver model, since SUMO does not take environmental conditions into account. There exists simulators that do count those factors, such as Carla \cite{dosovitskiy_carla_2017}who's built on Unreal engine, the backend is PhysX physics engine. Also metadrive \cite{li_metadrive_2022}, who built their system on bullet physics engine. However, one shared downside of them is that both simulators, although enabled by physics engine, they don't support scaling of scenes, for example, if i would like to evaluate the potentail outcome of a certain configuration of a workzone, I need to sequentially test all of them. This is requiring a lot of time. There's a gap in current AV siimulators, that that can both be scaled to multiple scene, multiple agents at one time, and also have phyiscs engine enabled.

% ------------------------------------------------------------------
%  P1 — Broad motivation: why simulation matters for AV
% ------------------------------------------------------------------

Simulators are indispensable to the development of autonomous driving.
Learning-based driving agents require platforms where their behavior can be evaluated in a controlled, reproducible, and safe manner—especially for rare or dangerous scenarios that cannot be ethically staged on public roads~\cite{grigorescu_av_survey_2020}.
There is a well-known gap between simulators and the real world, and researchers have been actively pursuing better simulation to narrow it.
A fundamental tension exists in simulator design: higher reconstruction fidelity—whether in physics, visuals, or both—demands greater computational cost, while systems that prioritize scalability, such as Waymax~\cite{gulino_waymax_2023} and GPUDrive~\cite{kazemkhani_gpudrive_2025}, simplify the underlying dynamics to bicycle models or other parametrized representations.
Recent GPU-parallel computing technologies have begun to shift this trade-off by moving heavy computation onto graphics processing units, opening the door to large-scale benchmarking for autonomous driving.
As Kalra and Paddock demonstrated, hundreds of millions to billions of miles of driving would be needed for a single autonomous vehicle to statistically demonstrate its reliability in terms of fatalities and injuries~\cite{kalra_driving_safety_2016}—a scale that is infeasible through real-world testing alone and underscores the necessity of high-throughput simulation.

% ------------------------------------------------------------------
%  P2 — Weather and safety motivation (from previous draft)
% ------------------------------------------------------------------

The safety relevance of simulation is further amplified by real-world driving conditions.
According to the Federal Highway Administration (FHWA)~\cite{noauthor_how_nodate}, approximately 6{,}000{,}000 vehicle crashes occur annually in the United States, of which roughly 12\% are weather-related.
Among those, 94\% involve changes in pavement friction: 73\% occur during rain or mist and 21\% during frozen precipitation.
Human driving behavior naturally adapts to adverse weather—drivers become more cautious when rainfall is high, compensating for reduced visibility and diminished maneuverability as effective road friction decreases~\cite{chen_influence_2019, ahmed_impacts_2018, peng_examining_2018}.
However, most existing simulation frameworks do not explicitly capture this behavioral adaptation, which may lead to unrealistic driving behavior in evaluation.
Accurate modeling of weather-dependent vehicle dynamics is particularly important for transportation planning, traffic safety analysis, and risk assessment of road infrastructure under variable conditions.

% ------------------------------------------------------------------
%  P3 — Traffic management decision-making motivation
% ------------------------------------------------------------------

Beyond autonomous-vehicle development, simulators also serve traffic administrators who make high-stakes operational decisions daily—road closures, lane reconfigurations, work zone designs—whose effects propagate into system-level economic and safety costs.
According to FHWA work zone operations guidelines, agencies actively monitor delay, queue length, and safety indicators, iteratively revising traffic management plans during execution~\cite{fhwa_workzone_2023}.
With the emergence of high-throughput simulators, administrators are increasingly moving from observing real system responses to evaluating candidate policies in simulation before deployment.
Yet, without accurate modeling of scene physics, ground friction, and plausible actor behavior, simulation-based decisions cannot be trusted, and suboptimal or dangerous policies may be selected.

% ------------------------------------------------------------------
%  P4 — Current simulators and their limitations
% ------------------------------------------------------------------

Traffic engineers have spent decades developing simulators at both microscopic and macroscopic levels.
SUMO~\cite{krajzewicz_sumo_2012}, for instance, is a widely used open-source microscopic traffic simulator employed in digital-twin applications, tracking individual vehicles with car-following and lane-changing models.
However, SUMO provides only limited support for weather effects and does not model tire--road friction physics.
Simulators such as CARLA~\cite{Dosovitskiy17}, built on Unreal Engine~4 with PhysX-based vehicle physics, and MetaDrive~\cite{li2022metadrive}, built on Panda3D with the Bullet engine, incorporate richer vehicle dynamics and environmental modeling.
However, in these systems each vehicle is driven by non-vectorized per-agent Python loops: actuation commands, state queries, and reward computations are executed sequentially for every agent, making the CPU-side control code---not the physics solver---the throughput bottleneck.
As a result, neither simulator supports \emph{GPU-parallel execution of multiple independent scenes}: evaluating $N$ scenarios requires $N$ sequential simulation runs, which is prohibitively slow for large-scale benchmarking or policy search.

On the other end of the spectrum, GPU-accelerated simulators such as Waymax~\cite{gulino_waymax_2023} and GPUDrive~\cite{kazemkhani_gpudrive_2025} achieve impressive throughput by batching hundreds of scenarios on a single GPU.
GPUDrive, built on the Madrona game engine, can simulate over a million steps per second for multi-agent training on the Waymo Open Motion Dataset.
However, these systems achieve scalability by simplifying the underlying dynamics. The vehicles follow bicycle or kinematic models with no rigid-body physics engine governing tire forces, suspension, or contact dynamics.
This simplification is reasonable for many research tasks, but limits applicability to safety-critical analyses where physical fidelity matters.

% ------------------------------------------------------------------
%  P5 — The gap statement
% ------------------------------------------------------------------

There exists a gap in the current landscape of autonomous-driving simulators: no existing system simultaneously provides
(i)~GPU-parallel execution of multiple independent driving scenes,
(ii)~multi-agent support within each scene,
(iii)~\emph{vectorized} articulated vehicles whose actuation, sensing, and reward computation scale with GPU parallelism rather than with per-agent CPU loops, and
(iv)~a full rigid-body physics engine governing vehicle dynamics with weather-aware friction modeling.

% ------------------------------------------------------------------
%  P6 — Our solution
% ------------------------------------------------------------------

To fill this gap, we present \textbf{SceneFactory}, a procedural scene construction and multi-agent training platform built on NVIDIA Isaac Sim and Isaac Lab~\cite{nvidia_isaac_2025}.
Isaac Lab provides a vectorized multi-agent reinforcement learning environment (DirectMARLEnv) that runs entirely on the GPU, with PhysX~5 as the underlying rigid-body physics engine.
A key technical contribution is a \emph{fully vectorized articulated PhysX vehicle}: all joint-state queries, actuation commands, observation assembly, and reward computation are expressed as batched tensor operations over the full agent population, eliminating the per-agent Python loops that throttle existing physics-based simulators.
This vectorized design achieves up to \textbf{127$\times$} higher throughput than a non-vectorized PhysX baseline on the same GPU and physics solver (\S\ref{sec:throughput}).
SceneFactory ingests real-world road topologies from the Waymo Open Motion Dataset~\cite{Kan_2024_icra}, procedurally constructs physically realistic USD road environments, and orchestrates the parallel simulation of hundreds of these articulated vehicles across dozens of distinct road worlds on a single GPU.
Additionally, we integrate a weather-to-friction module based on Zhao et al.'s tire--pavement friction model~\cite{zhao_tire-pavement_2024}, which maps precipitation conditions and road surface type to effective friction coefficients applied directly to PhysX material properties.
This paper presents the first phase of a larger research program; it describes the full pipeline and reports initial training results.

% ------------------------------------------------------------------
%  P7 — Contribution list
% ------------------------------------------------------------------

The contributions of this work are:
\begin{enumerate}[leftmargin=*, itemsep=2pt]
  \item \textbf{SceneFactory}, a procedural scene construction pipeline that converts Waymo Open Motion Dataset scenarios into GPU-parallelized, multi-world Isaac Sim stages with full PhysX rigid-body dynamics.
  \item A \textbf{fully vectorized articulated PhysX vehicle} whose joint-state queries, actuation, observation assembly, and reward computation are expressed entirely as batched GPU tensor operations, achieving up to 127$\times$ higher throughput than a non-vectorized PhysX baseline on the same hardware.
  \item A \textbf{multi-agent RL training environment} supporting $N$ parallel worlds $\times$ $M$ agents per world, all driven by a shared navigation policy trained via PPO.
  \item A \textbf{weather-to-friction module} that maps precipitation and road surface type to PhysX tire friction coefficients, enabling weather-aware driving behavior training.
\end{enumerate}

\chattodo{Insert concrete throughput and success-rate numbers once experiments are finalized (e.g., ``X env-steps/sec with Y agents on a single A6000 pro, Z\% goal-reach rate'').}

% ============================================================
%  sections/related_work.tex  —  Related Work (merged & cleaned)
%  Sources: SOC draft + sec/2_related_works.tex (previous writeup)
% ============================================================
\section{Related Work}
\label{sec:related_work}
%the next section we are going to introduce what kind urban driving simulators are already released, and we have a table to compare each of them, and by that authors can see that our simulator is the one that enables multiple scene, multiple agents, and have physics engine integrated, so that the agent response is more realistic. Ah forgot to say about it, but somewhere in the intro we might should say, it's reasonable of people to simplify some aspects of the system when building their simulators, but our simulator is more tailored to take care about safety related factors. I mean, surely most of the time, people drive in good weathers, especially waymo, or other autonomuos driving companies, their data collection is from those cities with good weathers, such as (idk idk, DC? California?) NOt sure how biased is their data, but the point is that cities like Madison, WI, or Buffalo, NY, those kinds of cities exists, and things like construction, communiting of civilians, and adverse weather, is unavoidably happening at the same time. In a simulator perspective, those are long tail events, but for traffic administrators, those are high stake scenarios in real life. Thus, simlator realism is one of the most cared aspect in this sense. Traffic administrators need to be able to compare the outcome of a certain decision, under a certain time of a day, and a certain kind of weather, to be able to make a  calculated decision. After introducing of those differnet current simulators, I should talk about what kind of studies in transportation exists about the friction between tyre and ground in different weather conditions. Since this is a work about driving simulator, the simulator related downstream tasks shall also be discussed in the realtaed work section. Inlcuding the use of NVIDIA cosmos transfer base model to generate visually realistic driving scenarios. 


% After the related work section, I need to actually write about methodology. 
% This work presents a GPU-accelerated, weather-aware, multi-world multi-agent training pipeline for physically grounded, safety-aware driving behavior.
% We review related work in three categories: driving simulators, weather-aware driving behavior research, and GPU-accelerated simulation frameworks.

% ==================================================================
\subsection{Driving Simulators}
\label{sec:rw_simulators}
% ==================================================================

As training robotic agents to drive in the real physical world is expensive and impractical, autonomous driving simulators are essential.
Multiple open-source platforms have been developed by the community, spanning a range of design trade-offs between physics fidelity, scene diversity, and computational scalability.

CARLA~\cite{Dosovitskiy17} is one of the earliest open-source autonomous driving simulators, built on the widely used Unreal Engine~4.
It is equipped with PhysX-based vehicle physics and provides multiple well-designed maps with extensive features and customization options.
However, its design makes it computationally expensive to run and unsuitable for GPU-accelerated workloads, limiting its applicability in tasks such as hardware-accelerated multi-world, multi-agent simulations.
Moreover, creating maps for CARLA is labor-intensive, as its focus is on highly detailed, handcrafted scenes rather than efficiently leveraging real-world data.

To address these limitations, MetaDrive~\cite{li2022metadrive} was introduced as another open-source simulator.
It aims to overcome limited scene diversity by allowing users to construct environments from open datasets such as the Waymo Open Motion Dataset~\cite{Kan_2024_icra}, or through procedural generation.
MetaDrive is backed by the Bullet physics engine (via Panda3D).
However, it does not support hardware acceleration, and only a single scene can be loaded at a time.
Although tire friction is reported as a configurable parameter in MetaDrive, its capability to train agents under varying friction conditions has not been demonstrated.

GPUDrive~\cite{kazemkhani_gpudrive_2025}, built on the Madrona game engine~\cite{shacklett_madrona_2023}, achieves impressive throughput for multi-scene, multi-agent training, capable of simulating over a million steps per second.
However, its vehicle dynamics use simplified collision and bicycle-model representations rather than a full rigid-body physics engine—a distinction that limits its ability to capture realistic tire--road interactions and contact dynamics.

A comparison of driving simulators across key capabilities is shown in Table~\ref{tab:simulators}.
Most physically realistic simulators lack hardware acceleration, and only a few incorporate realistic pavement friction modeling.
It is worth noting that simplification of dynamics is a valid design choice for many research tasks; SceneFactory targets the complementary use case of safety-critical evaluation where physics fidelity directly affects the validity of the analysis.

\begin{table*}[t]
\centering
\small
\begin{tabularx}{\textwidth}{l X X X X X X X X X}
\toprule
\textbf{Simulator} & \textbf{Multi-agent} & \textbf{GPU-Accel} & \textbf{Sensor Sim} & \textbf{Expert Data} & \textbf{Sim-agents} & \textbf{Physics Aware} & \textbf{Weather Aware}& \textbf{Routes / Goals} \\
\midrule
TORCS \cite{wymann_torcs_nodate}&  &  & \cmark &  & \cmark & \cmark &  & - \\
GTA V \cite{martinez_beyond_2017} &  &  & \cmark &  & & \cmark& \cmark & - \\
CARLA \cite{Dosovitskiy17} &  &  & \cmark &  & \cmark & \cmark & \cmark &Waypoints \\
Highway-env \cite{highway-env} &  &  &  &  &  & & & - \\
Sim4CV \cite{muller_sim4cv_2018} &  &  & \cmark &  &  & \cmark & &Directions \\
SUMMIT \cite{cai_summit_2020} & $\checkmark$ ($\geq 400$) &  & \cmark &  & \cmark & \cmark & \cmark & - \\
MACAD \cite{palanisamy_multi-agent_2019} & \cmark &  & \cmark &  & \cmark & \cmark & & Goal point \\
SMARTS \cite{zhou_smarts_2021} & \cmark &  &  &  &  & \cmark & & Waypoints \\
MADRaS \cite{santara_madras_2021} & $\checkmark$ ($\geq 10$) &  & \cmark &  &   & \cmark & & Goal point \\
DriverGym \cite{kothari_drivergym_2021} &  &  &  & \cmark & \cmark  & & & - \\
VISTA \cite{amini_vista_2022}& \cmark &  & \cmark & \cmark &   & & & - \\
nuPlan \cite{caesar_nuplan_2022} &  &  & \cmark & \cmark & \cmark &  & &Waypoints \\
Nocturne \cite{vinitsky_nocturne_2022} & $\checkmark$ ($\geq 128$) &  & \cmark & \cmark & \cmark & & & Goal point \\
MetaDrive \cite{li2022metadrive} & \cmark &  & \cmark & \cmark & \cmark & \cmark & \cmark & - \\
InterSim \cite{sun_intersim_2022} & \cmark &  &  & \cmark & \cmark &  & &Goal point \\
TorchDriveSim \cite{lavington_torchdriveenv_2024} & \cmark & \cmark &  &  &  &  & &- \\
BITS \cite{xu_bits_2022}& \cmark &  &  & \cmark & \cmark &  & &Goal point \\
Waymax \cite{gulino_waymax_2023} & $\checkmark$ ($\geq 128$) & \cmark &  & \cmark & \cmark &  & &Waypoints \\
GPUDrive \cite{kazemkhani_gpudrive_2025} & $\checkmark$ ($\geq 128$) & \cmark & \cmark & \cmark & \cmark &  & &Goal point \\
\textbf{SceneFactory (ours)} & $\checkmark$ ($\geq 128$) & \cmark & \cmark & \cmark & \cmark & \cmark & \cmark &Goal point \\
\bottomrule
\end{tabularx}
\caption{Comparison of driving simulators across key capabilities. SceneFactory is, to our knowledge, the first system that simultaneously provides GPU-accelerated multi-world execution, multi-agent support, full rigid-body physics, and weather-aware friction modeling.}
\label{tab:simulators}
\end{table*}

% ==================================================================
\subsection{Weather-Aware Driving Behavior}
\label{sec:rw_weather}
% ==================================================================

The relationship between weather conditions and road surface properties has been studied for decades in transportation research.
Adverse weather affects both driver visibility and vehicle dynamics: Chen et al.~\cite{chen_influence_2019} showed that drivers' perceived risk increases with worsening weather during car following, Ahmed and Ghasemzadeh~\cite{ahmed_impacts_2018} found that the probability of speed reduction increases by 23--29\% in rain using SHRP2 naturalistic data, and Peng et al.~\cite{peng_examining_2018} demonstrated significant changes in traffic parameters under fog and rain using real-time sensor data.

However, within the autonomous driving community, driving behavior under adverse weather has received comparatively less attention~\cite{gaonkar_enhancing_2025}.
Most autonomous driving systems have achieved strong performance under the implicit assumption of favorable weather~\cite{li_autonomous_2025}.
Existing research on adverse weather primarily focuses on perception and sensing—investigating how sensor performance degrades under rain, fog, or snow~\cite{song_weather_2020, li_autonomous_2025, fursa_worsening_2021, rahmati_edge_2025}—while comparatively fewer works explicitly consider how vehicle dynamics change.
Only recently have studies begun to explore driving behavior in low-friction environments using simulation; for example, Aghdasian et al.~\cite{aghdasian_tackling_2025} trained robust lane-keeping policies under snowy conditions in CARLA.

On the tire--pavement friction modeling side, Zhao et al.~\cite{zhao_tire-pavement_2024} developed a friction model that accounts for pavement texture and water film thickness, providing a physics-based mapping from environmental conditions to effective friction coefficients.
SceneFactory adopts this model to bridge the gap between weather conditions and vehicle dynamics within the simulation: precipitation and road surface type are converted to friction parameters that are applied directly to the PhysX physics engine, enabling agents to experience and learn from weather-dependent driving dynamics.

% ==================================================================
\subsection{GPU-Accelerated Simulation Frameworks}
\label{sec:rw_gpu_frameworks}
% ==================================================================

The need for large-scale RL training has driven the development of GPU-accelerated simulation engines.
Madrona~\cite{shacklett_madrona_2023} is a batch entity-component-system (ECS) engine that enables thousands of independent simulation worlds to run in parallel on a single GPU; it serves as the backend for GPUDrive.
Brax~\cite{freeman_brax_2021} provides a differentiable physics engine in JAX for rigid-body simulation, primarily targeting locomotion tasks.
MuJoCo~XLA extends the MuJoCo physics engine with JAX-based GPU parallelism but is not designed for vehicle-centric or USD-based scene management.

NVIDIA Isaac Sim and Isaac Lab~\cite{nvidia_isaac_2025, mittal_orbit_2023} take a different approach.
Isaac Sim provides a full simulation runtime built on the Omniverse platform with PhysX~5 as the rigid-body physics backend and OpenUSD as the scene description format.
Isaac Lab (evolved from Orbit) adds a vectorized environment framework on top, including the DirectMARLEnv abstraction for multi-agent RL with Fabric-based state cloning across parallel environments.
This combination—USD scene management, PhysX~5 articulated vehicle dynamics, and a GPU-resident training loop—makes Isaac Lab the natural choice for SceneFactory, where we require both high physics fidelity (articulated vehicles with tire friction, suspension, and contact sensing) and high throughput (hundreds of vehicles across dozens of parallel worlds).

\chattodo{Optionally add a brief paragraph on world models / Cosmos as a downstream application if the scope of the paper includes the video generation pipeline.}

% ============================================================
%  sections/simulator_design.tex  —  Simulator Design (Methodology)
% ============================================================
\section{Simulator Design}
\label{sec:simulator_design}

\chatstructure{This is the core technical section.  I recommend splitting it into the following subsections:
\begin{enumerate}
  \item \textbf{System Overview} — one-paragraph + architecture figure showing the full pipeline (Waymo $\to$ JSON $\to$ USD $\to$ Isaac Sim multi-world stage $\to$ RL training loop).
  \item \textbf{Data Processing \& Procedural Scene Construction} — the Waymo-to-USD pipeline, world layout, road prim generation (PointInstancer vs.\ explicit mesh), metadata baking.
  \item \textbf{Vehicle Asset Generation \& System Identification} — USD articulated vehicle, PhysX joint setup, CEM-based parameter fitting.
  \item \textbf{Weather-to-Friction Module} — Zhao et al.\ model, API design, PhysX material property updates.
  \item \textbf{Observation Space} — ego state, road geometry (KNN), neighbor vehicles, weather context.
  \item \textbf{Action Space} — throttle, steering, brake $\to$ PhysX joint targets.
  \item \textbf{Reward Design} — goal bonus, route progress, lane-keeping, collision/TTC penalties, off-road penalties.
  \item \textbf{Training Curriculum} — curriculum design (stages TBD).
\end{enumerate}
This mirrors the GPUDrive structure you mentioned, which is appropriate.}

As discussed in Section~\ref{sec:related_work}, physically realistic driving simulators such as CARLA and MetaDrive cannot execute multiple independent worlds in parallel, while batched simulators such as Waymax and GPUDrive simplify vehicle dynamics to bicycle or kinematic models.
SceneFactory bridges this gap.
We build on NVIDIA Isaac Sim and Isaac Lab~\cite{nvidia_isaac_2025, mittal_orbit_2023} and adopt elements of the observation and network design from GPUDrive~\cite{kazemkhani_gpudrive_2025}, specifically the $k$-nearest-neighbor road-point encoding and per-entity MLP with max-pool aggregation.
Our new contributions include procedural USD scene construction from Waymo data, articulated PhysX vehicle dynamics with system identification, and a physics-based weather-to-friction module.
This section describes each component.

% ==================================================================
\subsection{System Overview}
\label{sec:system_overview}
% ==================================================================

\paragraph{Platform choice.}
SceneFactory requires a simulation backend that can run a full physics engine across many independent worlds simultaneously.
NVIDIA Isaac Sim meets both requirements.
It is built on the Omniverse platform with PhysX~5 as the dynamics solver and OpenUSD as the scene description format, allowing vehicles to be modeled as articulated rigid bodies with suspension, steering, and wheel-spin joints.
Isaac Lab~\cite{mittal_orbit_2023} extends Isaac Sim with a vectorized reinforcement learning environment layer, including the \texttt{DirectMARLEnv} abstraction for multi-agent tasks.
Isaac Lab's cloning mechanism replicates a template scene across $N$ parallel environments on the GPU.
Its Fabric-based state management keeps the entire training loop (physics stepping, observation computation, reward evaluation, and policy inference) on-device, eliminating host-device data transfers in the inner loop.

\paragraph{Pipeline overview.}
The SceneFactory pipeline consists of four stages, illustrated in Figure~\ref{fig:architecture}:

\begin{enumerate}[leftmargin=*, itemsep=2pt]
  \item \textbf{Data ingestion.}
  Raw Waymo Open Motion Dataset~\cite{Kan_2024_icra} TFRecord files are converted offline into per-scenario JSON files.
  Each JSON encodes the road topology (typed polylines for lane centers, road edges, and road lines) and agent metadata (vehicle start poses, goal positions, and trajectory data).

  \item \textbf{Procedural scene construction.}
  A scene builder reads the JSON manifests and constructs USD road geometry for each scenario.
  Multiple scenarios are arranged in a 2D grid layout (configurable spacing, typically 200\,m between world centers) to form a single composite stage.
  Road geometry metadata (polyline coordinates, directions, lane types, and half-widths) is baked into USD \texttt{customData} attributes on each world's root prim for constant-time runtime access.
  Vehicle start and goal positions are extracted, centered, and filtered per scenario.

  \item \textbf{Environment instantiation.}
  Isaac Lab clones the template environment $N$ times.
  After cloning, SceneFactory builds road geometry independently inside each environment, so every parallel world can host a different Waymo-derived scenario.
  Procedurally generated vehicle assets (four-wheel suspension, front-wheel steering, front-wheel drive) are spawned at scenario-specified start poses, with up to $M$ agents per world.
  The weather-to-friction module (Section~\ref{sec:weather_friction}) computes per-world friction coefficients that are encoded into the policy's observation; per-world physics material variation is planned but not yet active (see Section~\ref{sec:weather_friction} for details).



  \item \textbf{Training loop.}
  The multi-agent environment is wrapped into a shared-policy vectorized interface that exposes one rollout slot per agent, yielding $N \times M$ total slots.
  A PPO policy~\cite{schulman_proximal_2017}, trained via the RSL-RL framework~\cite{schwarke_rsl-rl_2025}, consumes structured observations (ego state, road geometry context, neighbor vehicles, and optional weather context) and produces continuous throttle, steering, and brake commands decoded into joint targets.
  All inference, advantage estimation, and gradient updates remain on-device.
\end{enumerate}

\chattodo{\textbf{ACTION ITEM — FIGURE:} Create the architecture / pipeline diagram (Figure~\ref{fig:architecture}).  Suggested content: Waymo TFRecord $\to$ JSON $\to$ USD multi-world stage $\to$ Isaac Lab DirectMARLEnv $\to$ PPO training via RSL-RL.  This is arguably the most important figure in the paper.  Consider a horizontal flow with callout boxes for the weather-to-friction module and the vehicle sysid module feeding into the environment instantiation stage.}
\chattodo{\textbf{ACTION ITEM — LABEL:} Define \texttt{fig:architecture} once the figure file is added to \texttt{figures/}.}

\paragraph{Design principles.}
Three principles guided SceneFactory's design.
First, \emph{physics fidelity}: vehicles are articulated bodies with suspension, tire forces, and contact events resolved by the physics engine, rather than kinematic proxies.
Second, \emph{scene diversity}: any Waymo scenario can be loaded without recompilation; adding a scenario requires only a JSON file and a configuration entry.
Third, \emph{throughput}: the entire training loop stays on-device, allowing a shared navigation policy to be trained across thousands of agents in diverse road environments on a single GPU.
The following subsections detail each component.

\subsection{Data Processing and Procedural Scene Construction}
\label{sec:scene_construction}
\chattodo{Describe: (1) Waymo TFRecord to JSON conversion — what fields are extracted (road polylines, agent trajectories, start-goal pairs); (2) world layout — grid arrangement with 200m spacing; (3) road prim generation — PointInstancer mode vs.\ explicit-prim mode; (4) metadata baking into USD custom data attributes for constant-time runtime access.  This is a genuine engineering contribution — GPUDrive compiles scenarios into C++ at build time, your approach via USD is more flexible.}
\chatcite{Ettinger et al., 2021 (Waymo Open Motion Dataset).}


SceneFactory consumes driving scenarios from the Waymo Open Motion Dataset~\cite{Kan_2024_icra} and constructs physically simulated USD environments through a three-step offline pipeline: data extraction, per-scenario world building, and multi-world grid assembly.

\paragraph{Waymo TFRecord to JSON.}
Each Waymo TFRecord contains up to 30{,}000 roadgraph sample points and 128 tracked agents across 91 timesteps.
An offline conversion script extracts two categories of information per scenario.
First, \emph{road topology}: sample points are grouped by polyline ID, and points within each polyline are PCA-ordered along the dominant XY axis to ensure consistent traversal direction.
Each polyline retains its Waymo-assigned type code (e.g., lane centers are types 1--2, road edges are types 15--16), per-point 3D coordinates, and per-point direction vectors.
Second, \emph{agent metadata}: for each valid vehicle track, the script records the start pose (position and yaw at the first observed timestep) and the goal pose (position and yaw at the last observed timestep).
The output is one JSON file per scenario, containing a \texttt{road.polylines} array and an \texttt{agents.items} array.

\paragraph{Per-scenario world building.}
Given a scene JSON, a world builder constructs USD road geometry through the following steps:
\begin{enumerate}[leftmargin=*, itemsep=1pt]
  \item \textbf{Re-centering.}
  A scene center is computed as the mean of all road polyline points.
  All coordinates (road and agent) are shifted so that each scenario is locally centered at the origin.

  \item \textbf{Z-flattening.}
  Waymo road elevation data is collapsed to a constant ground height ($z=0$). Currently this is a simplified design, and future work will extend the driving to include elevation of the terrain as well.

  \item \textbf{Segment generation.}
  Each consecutive pair of polyline points defines an oriented box segment.
  The segment is placed at the midpoint, rotated to align with the point-to-point direction, and assigned a configurable width (default 0.10\,m) and height (default 0.10\,m).
  Segment pairs separated by more than 3.0\,m are treated as discontinuities and skipped.
  Both endpoints must fall within a bounding box of $\pm B/2$ from the scene center ($B=200$\,m in our experiments) to be included.

  \item \textbf{USD prim creation.}
  Road segments are grouped by polyline type.
  Each group is rendered as a \texttt{PointInstancer} with a single unit-cube prototype; per-segment position, orientation, and scale arrays encode the full geometry compactly.
  Each type group receives a distinct material color for visualization.

  \item \textbf{Metadata baking.}
  After all segments are created, five arrays are written to the world root prim's USD \texttt{customData}: segment midpoint positions, direction vectors, polyline type codes, half-lengths (half the point-to-point distance), and half-widths.
  These arrays encode the full road geometry in a flat, GPU-friendly format that can be loaded at runtime without re-parsing the scene JSON or traversing the USD prim hierarchy.
\end{enumerate}

\paragraph{Agent spawn extraction.}
Vehicle start and goal positions are parsed from the same scene JSON.
Coordinates are re-centered using the same scene center as the road geometry.
Agents whose start position falls outside the bounding box are discarded.
Agents whose start position is already within the goal radius (default 3.0\,m) are optionally filtered out to avoid trivial episodes.
The remaining agents are capped at $M$ per world (e.g., $M{=}16$ in training).

\paragraph{Multi-world grid assembly.}
To populate $N$ parallel environments, SceneFactory arranges worlds in a 2D grid.
Each world occupies a cell of size $200 \times 200$\,m with configurable padding (typically 200\,m), yielding a pitch of 400\,m between world centers.
The world root prim receives an XformOp translation to place it at the correct grid position.
World-to-environment assignment is controlled by the curriculum configuration: in \emph{fixed} mode all environments load the same scenario, while in \emph{random-fill} mode a curated list of scenarios is shuffled and tiled to fill $N$ slots.

\paragraph{Post-clone road injection.}
Isaac Lab's environment cloning mechanism replicates a single template environment $N$ times.
However, SceneFactory requires each environment to host a \emph{different} road scenario.
We solve this by building road geometry \emph{after} Isaac Lab's clone step: the template environment contains only the vehicle articulations and ground plane, and per-environment road prims are injected into the already-cloned USD stage.
This avoids cloning complex instancer prims (which is unreliable under Fabric) and allows heterogeneous road layouts across all environments.

\paragraph{GPU metadata tensors.}
At simulation startup, the baked \texttt{customData} arrays are read from each environment's world root prim and loaded into padded GPU tensors of shape $(N, P_{\max}, d)$, where $P_{\max}$ is the maximum segment count across all environments and $d$ is the feature dimension.
A boolean validity mask distinguishes real segments from padding.
These tensors support vectorized lane-proximity queries and oriented-box overlap tests at runtime, enabling reward computation and observation encoding without any USD access during training.


\chattodo{\textbf{ACTION ITEM --- FIGURE:} Consider adding a small figure showing the JSON-to-USD pipeline for one scenario: raw Waymo polylines $\to$ re-centered segments $\to$ USD instancer prims $\to$ multi-world grid.}


\subsection{Vehicle Asset Generation and System Identification}
\label{sec:vehicle_sysid}
\chattodo{Describe: (1) how the articulated PhysX vehicle USD asset is constructed (4 wheels, front-wheel steering revolute joints, rear-wheel drive effort joints, contact sensors); (2) CEM-based system identification to fit PhysX parameters (tire friction, suspension damping/stiffness) against reference Waymo trajectory data; (3) report the tracking error from sysid — this quantifies what ``more realistic'' actually means.}
\chatcomment{The sysid results are critical.  If you claim physics fidelity as a differentiator over GPUDrive, you need numbers showing that your PhysX vehicle tracks reference trajectories better than a bicycle/kinematic model.}

SceneFactory uses a procedurally generated articulated vehicle rather than a hand-authored asset.
The vehicle is defined by a parametric specification, converted to URDF, imported into Isaac Sim as a USD asset, and then calibrated via system identification.

\paragraph{Articulated vehicle structure.}
The vehicle is a 10-DOF articulated rigid body.
The kinematic chain is rooted at a box chassis ($4.0 \times 2.0 \times 1.0$\,m, 1800\,kg) and branches to four corners, each with the following joint chain:
%
\begin{itemize}[leftmargin=*, itemsep=1pt]
  \item A \textbf{prismatic} joint (Z-axis) models the suspension, with travel $\pm0.175$\,m.
  \item On the two front corners only, a \textbf{revolute} joint (Z-axis) models steering, with limits $\pm32^{\circ}$.
  \item A \textbf{continuous} joint (Y-axis) models wheel spin.
\end{itemize}
%
Each wheel is a cylinder (radius 0.35\,m, width 0.15\,m, 20\,kg).
The wheelbase is 2.6\,m and the track width is 2.0\,m.
Inertia tensors for all links are computed from solid-body formulas (box for chassis and suspension links, cylinder for wheels).
The URDF is generated programmatically from a frozen dataclass (\texttt{StudentVehicleSpec}) and imported into Isaac Sim via the URDF importer with joint-drive and self-collision settings enabled.
Post-import processing authors PhysX contact materials (static and dynamic friction $= 1.0$, restitution $= 0.0$, compliant contact) and de-instances mesh prims for compatibility with Fabric.

\paragraph{Action decoding.}
The policy outputs three continuous actions per agent: throttle $\in [0,1]$, steering $\in [-1,1]$, and brake $\in [0,1]$.
These are decoded into PhysX joint commands as follows.
Steering is converted to a position target $\theta_{\mathrm{target}} = \text{steering} \times \theta_{\max}$ and applied via a PD controller:
\begin{equation}
\label{eq:steer_pd}
\tau_{\mathrm{steer}} = \mathrm{clamp}\!\bigl(K_p\,(\theta_{\mathrm{target}} - \theta) - K_d\,\dot{\theta},\;\pm\tau_{\mathrm{steer,max}}\bigr).
\end{equation}
Throttle applies a drive torque to the two front wheels (front-wheel drive): $\tau_{\mathrm{drive}} = \text{throttle} \times \tau_{\mathrm{drive,max}}$.
Brake applies an opposing torque to all four wheels, with sign determined by wheel angular velocity (and a sign-memory mechanism to handle zero-velocity transitions).
An external wrench on the chassis provides two stabilization terms: a lateral velocity damping force $F_{\mathrm{lat}} = -\lambda_{\mathrm{lat}} \cdot f_{\mathrm{lat}} \cdot v_{\mathrm{lat}}$ and a yaw damping torque $\tau_z = -\lambda_{\mathrm{yaw}} \cdot f_{\mathrm{lat}} \cdot \omega_z$, where $f_{\mathrm{lat}}$ is the sysid-tuned lateral friction scale for the active surface type.
The concrete PD gains and torque limits are listed in Table~\ref{tab:action_params} (\S\ref{sec:actions}).

% \chattodo{Report the final PD gains ($K_p$, $K_d$) and torque limits from the sysid result used in experiments.}

\paragraph{System identification.}
The vehicle model has 20 tunable dynamics parameters: drive, brake (front/rear), and steering torques; steering PD gains and limits; suspension stiffness and damping; wheel mass and inertia scale; chassis center-of-mass offset; lateral and yaw damping coefficients; and per-surface friction scales for three surface types (dry, wet, gravel).
We fit these parameters using the cross-entropy method (CEM).

\emph{Reference data.}
A PhysX ``teacher'' vehicle (the Isaac Sim built-in PhysX vehicle model, not our articulated URDF vehicle) executes seven scripted maneuvers on a multi-surface patch track: straight-line acceleration and braking, left and right step-steer inputs, constant-radius turns, sinusoidal steering, and an S-curve surface transition across dry asphalt ($\mu_s{=}1.0$), wet asphalt ($\mu_s{=}0.75$), and gravel ($\mu_s{=}0.60$).
Each rollout records per-frame chassis position, yaw, speed, yaw rate, wheel angular velocities, and steer angles at 60\,Hz.

\emph{Loss function.}
Let $x^{\mathrm{ref}}_{k,t}$ denote the teacher's value for channel $k$ at frame $t$, and $x^{\mathrm{stu}}_{k,t}$ the corresponding student value under a candidate parameter vector.
Each rollout has $T$ frames recorded at 60\,Hz.
The per-channel mean squared error is $\mathrm{MSE}_k = \tfrac{1}{T}\sum_{t=1}^{T}(x^{\mathrm{stu}}_{k,t} - x^{\mathrm{ref}}_{k,t})^2$.
The full objective adds two terminal-frame penalties that emphasize end-of-maneuver accuracy:
\begin{equation}
\label{eq:sysid_loss}
\mathcal{L} = \sum_{k \in \mathcal{K}} w_k \cdot \mathrm{MSE}_k
  + w_{\mathrm{pos,final}} \cdot \|\mathbf{p}^{\mathrm{stu}}_T - \mathbf{p}^{\mathrm{ref}}_T\|^2
  + w_{\mathrm{spd,final}} \cdot (s^{\mathrm{stu}}_T - s^{\mathrm{ref}}_T)^2,
\end{equation}
where $\mathcal{K} = \{\text{position}_{xy}, \text{yaw}, \text{speed}, \text{yaw rate}, \text{wheel speed}, \text{steer angle}\}$ with weights $w_k = (1.0, 0.4, 0.35, 0.25, 0.05, 0.05)$ respectively.
The vectors $\mathbf{p}^{\mathrm{stu}}_T, \mathbf{p}^{\mathrm{ref}}_T \in \mathbb{R}^2$ are the student and teacher XY positions at the final frame, and $s^{\mathrm{stu}}_T, s^{\mathrm{ref}}_T$ are the corresponding speeds.
The terminal weights are $w_{\mathrm{pos,final}} = 1.5$ and $w_{\mathrm{spd,final}} = 0.75$.


\emph{Multi-stage search.}
The CEM search is divided into five stages to manage the 20-dimensional space:
(1)~\emph{longitudinal}: tunes drive/brake torques and wheel parameters on the straight-line rollout;
(2)~\emph{steering}: tunes steering gains, suspension, and damping on the turning rollouts;
(3)~\emph{surface}: tunes per-surface friction scales on the transition rollout;
(4)~\emph{refinement}: tunes all 20 parameters jointly on all rollouts with a narrowed search window (18\% of the full range);
(5)~\emph{brake preservation}: re-tunes longitudinal parameters with a tight window (10\%) to prevent regression.
Each stage inherits the best configuration from the previous stage.
The CEM uses a population of 16 candidates, an elite fraction of 25\%, and an initial standard deviation of 25\% of each parameter's range.

\chattodo{\textbf{ACTION ITEM --- TABLE:} Add a table reporting per-stage sysid loss (total, position MSE, speed MSE, yaw MSE, yaw-rate MSE) from the final sysid run.  The numbers from the \texttt{fwd\_v1\_staged\_cem\_anchor\_overnight} run are: Stage~1 total=0.39, Stage~4 (refinement) total=9.35.  Consider also reporting a comparison against the default (untuned) vehicle.}
\chattodo{\textbf{ACTION ITEM --- FIGURE:} Consider a trajectory overlay figure showing teacher vs.\ student vehicle on one or two maneuvers (e.g., step-steer left and surface transition).}

\subsection{Weather-to-Friction Module}
\label{sec:weather_friction}
\chattodo{Describe the Zhao et al.\ (2024) friction model: inputs (water film thickness $h$, road surface type), output (effective friction coefficient $\mu$).  Show the API: user sets precipitation level $\to$ module computes $h$ $\to$ Zhao model gives $\mu$ $\to$ $\mu$ written to PhysX material properties per world.  Include the formula or at least a qualitative plot of $\mu$ vs.\ $h$ for different surface types.}


Wet-road friction is a first-order factor in vehicle dynamics, yet most driving simulators either ignore it or expose only a single user-specified friction scalar.
SceneFactory integrates a physics-based tire-pavement friction model that maps weather conditions to per-world friction coefficients, enabling the policy to train across a realistic range of road-surface conditions.

\paragraph{Friction model.}
We implement the modified Average Lumped LuGre (ALL) model of Zhao et al.~\cite{zhao_tire-pavement_2024}, which extends the LuGre bristle-deflection framework to account for pavement texture and water film thickness.
The model evolves the average bristle deflection $\bar{z}$ via the ODE
\begin{equation}
\label{eq:bristle_ode}
\frac{d\bar{z}}{dt} = v_r - \theta \, Y_R \left[\frac{\sigma_0 \lvert v_r \rvert}{g(v_r)}\,\bar{z} + K \lvert \omega_r \rvert\,\bar{z}\right],
\end{equation}
where $v_r$ is the relative (slip) speed between the tire contact patch and the road, $\omega_r$ is the wheel circumferential speed, $\sigma_0$ is the bristle stiffness, $K = 7/(6L)$ is the ALL approximation constant with $L$ the contact patch length, $\theta$ is a texture influence coefficient (per road surface type), and $Y_R \in [0,1]$ is the contact patch length ratio under hydrodynamic lift.
The Stribeck steady-state function $g(v_r) = \mu_c + (\mu_s - \mu_c)\exp\bigl(-\lvert v_r / v_s \rvert^\alpha\bigr)$ captures the transition from static friction $\mu_s$ to Coulomb friction $\mu_c$.
The Stribeck speed depends on the water film thickness $h_w$ (in meters) via
\begin{equation}
\label{eq:stribeck_speed}
v_s(h_w) = b_1 \exp(1000\,b_2\,h_w + b_3) + b_4,
\end{equation}
where $b_1 = 4.8916$, $b_2 = -7.91$, $b_3 = 3.01$, and $b_4 = 3.40$ are calibrated constants from~\cite{zhao_tire-pavement_2024}.
Because $b_2 < 0$, $v_s$ decreases with water film thickness: thicker water films cause the friction drop-off to onset at lower slip speeds.

Given the steady-state solution $\bar{z}^*$, the effective friction coefficient is
\begin{equation}
\label{eq:mu_all}
\mu = \max\!\bigl(\theta\,Y_R - Y_F,\; 0\bigr) \cdot \bigl(\theta\,Y_R\,\sigma_0\,\bar{z}^* + \sigma_1\,\dot{\bar{z}} + \sigma_2\,v_r\bigr),
\end{equation}
where $Y_F \in [0,1]$ is the hydrodynamic force ratio and $\sigma_1$, $\sigma_2$ are viscous damping coefficients.
Under the default calibration from~\cite{zhao_tire-pavement_2024}, both $\sigma_1 = 0$ and $\sigma_2 = 0$, so the effective formula reduces to $\mu = \max(\theta\,Y_R - Y_F,\,0)\cdot\theta\,Y_R\,\sigma_0\,\bar{z}^*$.
Both $Y_R$ and $Y_F$ depend on vehicle speed, water film thickness, tire geometry (contact patch dimensions, tread radius, wheel radius), and physical constants (water density and viscosity).
When the contact factor $\theta\,Y_R - Y_F$ drops to zero, the model predicts full hydroplaning.
Our implementation numerically integrates Equation~\ref{eq:bristle_ode} to steady state using exponential integration ($\Delta t = 10^{-3}$\,s, convergence threshold $\epsilon = 10^{-7}$), matching the paper's calibrated parameters (Table~4 and Table~5 in~\cite{zhao_tire-pavement_2024}).

\paragraph{Road surface presets.}
The texture influence coefficient $\theta$ and the texture amplitude $T_d$ vary by pavement type.
We provide three presets calibrated to the values in~\cite{zhao_tire-pavement_2024}:
\emph{Asphalt Concrete} (AC, $\theta{=}1.00$, $T_d{=}0.65$\,mm),
\emph{Stone Mastic Asphalt} (SMA, $\theta{=}1.09$, $T_d{=}0.80$\,mm), and
\emph{Open-Graded Friction Course} (OGFC, $\theta{=}1.21$, $T_d{=}1.08$\,mm).
OGFC, designed for drainage, retains higher friction under wet conditions; AC, the most common pavement, serves as the baseline.

\paragraph{Pipeline.}
Each world in the training stage is assigned a friction configuration specifying the road surface type and the water film thickness $h_w$ (in mm).
At environment construction time, the pipeline feeds these values into the modified ALL model (Equation~\ref{eq:mu_all}) at a reference speed of 13.89\,m/s ($\approx$50\,km/h) and two slip ratios (0.15 for static, 0.80 for dynamic) to produce per-world static and dynamic friction coefficients $\mu_s$ and $\mu_d$.
The effective coefficient $\mu_{\mathrm{eff}}$ (static by default) is used downstream as follows.

The ground-plane material uses the sysid-tuned friction scale for the dry-asphalt surface type (Section~\ref{sec:vehicle_sysid}), with static friction $\mu^{\mathrm{ground}}_s = \min(1.0,\; f_{\mathrm{surface}} \cdot \sqrt{f_{\mathrm{lon}} \cdot f_{\mathrm{lat}}})$ and dynamic friction $\mu^{\mathrm{ground}}_d = 0.95\,\mu^{\mathrm{ground}}_s$.
PhysX resolves contact friction using the minimum of the ground and tire material values (combine mode \texttt{min}).
Because all cloned environments share a single global ground plane, this material is uniform across worlds.
Per-world friction variation is currently communicated to the policy only through the observation: the weather conditions are encoded into a 4-dimensional vector appended to the ego state:
\begin{equation}
\label{eq:weather_obs}
\mathbf{o}_{\mathrm{weather}} = \bigl[\, h_w / h_{\mathrm{norm}},\; \mathbb{1}_{\mathrm{AC}},\; \mathbb{1}_{\mathrm{SMA}},\; \mathbb{1}_{\mathrm{OGFC}} \,\bigr] \in \mathbb{R}^4,
\end{equation}
where $h_{\mathrm{norm}} = 1.0$\,mm is a normalization constant and $\mathbb{1}_{\{\cdot\}}$ is a one-hot encoding of the road surface type.
This allows the policy to condition its driving behavior on surface conditions, but the physics solver does not yet enforce different contact friction across worlds.

\chattodo{\textbf{ACTION ITEM --- PER-WORLD FRICTION:} The current implementation does not apply per-world friction to the physics.  The ground plane is a single global prim shared across all environments, making per-world ground materials infeasible under Isaac Lab's cloning.  A viable alternative is to assign per-vehicle wheel materials at spawn time: since each vehicle's wheel collision meshes live inside their own \texttt{/World/envs/env\_X/} scope, their PhysX contact materials can be set independently.  Under the \texttt{min} friction combine mode, lowering the wheel material friction to $\mu_{\mathrm{eff}}$ would produce the correct tire-ground friction.  Implement this before camera-ready.}

\chattodo{\textbf{ACTION ITEM --- FIGURE:} Consider a plot of $\mu_{\mathrm{eff}}$ vs.\ $h_w$ (0--1.0\,mm) for the three surface presets (AC, SMA, OGFC) at $v{=}$13.89\,m/s.  The \texttt{python -m src.trfc --demo} CLI produces a sweep that can generate this data.}

\chatcite{Zhao et al., 2024 \checkmark{}.}

\subsection{Observation Space and Network Architecture}
\label{sec:observations}
\chattodo{Detail the 1,929-dim observation: ego state (7-dim), road geometry context (350 points $\times$ 5-dim = 1,750-dim), neighbor vehicles (24 $\times$ 7-dim = 168-dim), weather context (4-dim).  Describe the late-fusion architecture: per-group encoder MLPs, max-pool aggregation for variable-size sets, fusion trunk.}
\chatcomment{A figure showing the observation encoding and network architecture would be very helpful here.}

Each agent receives a 1{,}929-dimensional observation vector assembled from four groups: ego state, weather context, road geometry context, and neighbor vehicles.
All spatial quantities are expressed in the agent's body frame via a 2D rotation by $-\psi$ (ego heading), so the policy sees an ego-centric view of the world.
The groups are processed by a late-fusion actor-critic network with per-modality encoders and max-pool set aggregation, following the design of GPUDrive~\cite{kazemkhani_gpudrive_2025}.

\paragraph{Ego state (7-dim).}
The ego state encodes the agent's goal and velocity in body-frame coordinates:
%
\begin{enumerate}[leftmargin=*, itemsep=1pt]
  \item Goal position $(g_x^b, g_y^b)$, each divided by the scene bounding-box half-size ($B/2 = 100$\,m).
  \item Heading error to the goal, encoded as $(\sin\Delta\psi,\,\cos\Delta\psi)$.
  \item Euclidean goal distance $\|g^b_{xy}\|$, divided by $B/2$.
  \item Body-frame velocity $(v_x^b, v_y^b)$, each divided by 10\,m/s.
\end{enumerate}
%
When the weather-to-friction module is active, the 4-dim weather context (Equation~\ref{eq:weather_obs}) is appended to the ego group, yielding an 11-dim input to the ego encoder.

\paragraph{Road geometry context ($K_r \times 5 = 1{,}750$-dim).}
Each agent selects $K_r = 350$ road-segment midpoints from the pre-loaded GPU metadata tensors (Section~\ref{sec:scene_construction}).
The selection operates in two steps: first, all road points within a radius $r = 10$\,m of the ego position are identified via a squared-distance threshold; second, these candidate points are sorted by their original polyline index to preserve road topology, and the first $K_r$ are taken.
If fewer than $K_r$ points fall within the radius, the remaining slots are zero-padded.
Each selected point contributes five features:
%
\begin{enumerate}[leftmargin=*, itemsep=1pt]
  \item Relative position $(\Delta x^b, \Delta y^b)$ in ego body frame, each divided by $r$.
  \item Polyline type code, divided by a normalizer ($= 20$).
  \item Road direction unit vector $(d_x^b, d_y^b)$ in ego body frame.
\end{enumerate}

\paragraph{Neighbor vehicles ($K_v \times 7 = 168$-dim).}
The $K_v = 24$ nearest alive agents (by Euclidean distance in world XY) are selected per ego agent.
Self-distance and dead-agent distances are set to infinity before the argsort.
Each neighbor contributes seven features:
%
\begin{enumerate}[leftmargin=*, itemsep=1pt]
  \item Relative position $(\Delta x^b, \Delta y^b)$ in ego body frame, each divided by $B/2$.
  \item Chassis length and width, each divided by $B/2$.
  \item Relative heading $\Delta\psi$, divided by $\pi$.
  \item Speed $\|v_{xy}\|$, divided by 10\,m/s.
  \item Pairwise time-to-collision (TTC), computed via a swept-circle test.
    Each vehicle is approximated by three equal-radius circles placed along the longitudinal axis at offsets $\{-d,\, 0,\, +d\}$ from the chassis centre, where $r = \max(0.45,\; 0.55\, W)$ and $d = \min\!\bigl(\tfrac{\mathit{WB}}{2},\; \max(0,\; \tfrac{L}{2} - 0.8\,r)\bigr)$.
    For the default vehicle ($L{=}4.0$\,m, $W{=}2.0$\,m, $\mathit{WB}{=}2.6$\,m) this gives $r \approx 1.10$\,m and $d \approx 1.12$\,m.
    The test solves a quadratic closest-approach equation for every pair among the $3 \times 3 = 9$ circle combinations and takes the minimum.
    The resulting TTC is clamped to $[0,\, \tau_{\max}]$ and divided by $\tau_{\max} = 10$\,s.
    Neighbors behind the ego or with no collision trajectory receive $\mathrm{TTC}/\tau_{\max} = 1$.
\end{enumerate}
%
Slots beyond the number of alive neighbors are zero-padded.

\paragraph{Late-fusion actor-critic.}
We adopt the per-entity encoder with max-pool aggregation architecture introduced by GPUDrive~\citep{kazemkhani_gpudrive_2025}, with two modifications: (i)~we use \emph{masked} max-pool (zero-padded slots are filled with $-\infty$ before the $\max$) rather than a plain pool, and (ii)~we replace their Tanh activations with ELU.
The flat observation is split into its three groups and processed by separate encoder MLPs, one set for the actor and one for the critic (no weight sharing).
%
%
\begin{itemize}[leftmargin=*, itemsep=1pt]
  \item \textbf{Ego encoder}: $11 \to 64 \to 64$ (two hidden layers, ELU activations).
  \item \textbf{Road encoder}: $5 \to 96 \to 96$, applied independently to each of the $K_r$ points; the $K_r$ output vectors are aggregated via \emph{masked max-pool}, where zero-padded points are filled with $-\infty$ before the $\max$ operation.
  \item \textbf{Vehicle encoder}: $7 \to 96 \to 96$, applied independently to each of the $K_v$ neighbors; aggregated via masked max-pool in the same way.
\end{itemize}
%
The three embeddings (64 + 96 + 96 = 256) are concatenated and passed through a shared trunk MLP ($256 \to 128 \to 64$, ELU).
The actor head projects to a 3-dim mean vector $\boldsymbol{\mu}$ (one per action), paired with a state-independent learnable log-standard-deviation $\log\boldsymbol{\sigma}$ to form a diagonal Gaussian policy.
The critic head projects to a scalar state-value estimate.
No empirical observation normalization (running mean/std) is used; all features are manually scaled at assembly time as described above.

\chattodo{\textbf{ACTION ITEM --- FIGURE:} Create a network architecture diagram showing the three encoder branches, max-pool aggregation, fusion trunk, and actor/critic heads.  This is the second most important figure after the pipeline diagram.}
\chatcite{Kazemkhani et al., 2025 \checkmark{} — for the KNN road-point encoding and per-entity MLP + max-pool design.}


\subsection{Action Space}
\label{sec:actions}
\chattodo{Three continuous actions: throttle $\in [0,1]$, steering $\in [-1,1]$, brake $\in [0,1]$.  Decoded into PhysX joint position targets (steering) and effort targets (drive/brake).  Keep this brief.}


Each agent outputs a three-dimensional continuous action $\mathbf{a} = (a_{\mathrm{thr}},\, a_{\mathrm{str}},\, a_{\mathrm{brk}})$.
The policy produces values in $[-1,1]^3$; throttle and brake are rectified to $[0,1]$ before decoding, so a zero-mean Gaussian policy naturally idles.
All three channels are decoded into PhysX \emph{effort targets} (no position or velocity targets are used at runtime); the steering PD loop described in Eq.\ref{eq:steer_pd} is computed in software at every physics substep.

\paragraph{Control frequency.}
Physics runs at $1/\Delta t_{\mathrm{phys}} = 120$\,Hz.
Actions are held constant over a decimation window of 4 substeps, yielding a control frequency of 30\,Hz.
Within each substep the steering PD reads the current joint state, so the effective PD loop rate is 120\,Hz.

\paragraph{Brake sign memory.}
Brake torque must oppose wheel rotation, so its sign is $-\mathrm{sign}(\dot{q}_{\mathrm{wheel}})$.
When $|\dot{q}_{\mathrm{wheel}}| < 10^{-4}$\,rad/s the sign becomes unreliable; we therefore latch the last observed sign and continue applying it until rotation resumes.
This prevents torque chatter near zero velocity.

\paragraph{Actuation parameters.}
Table~\ref{tab:action_params} lists the sysid-tuned values used in all experiments.
Drive torque is applied to the two front wheels only (FWD); brake torque is applied to all four wheels with separate front/rear magnitudes.
Both drive and brake torques are further scaled by a shared longitudinal scale $s_{\mathrm{long}} = 0.9273$, which is one of the 20 sysid parameters (\S\ref{sec:vehicle_sysid}).

\begin{table}[h]
  \centering\small
  \caption{Action-decoding parameters (sysid-tuned values).}
  \label{tab:action_params}
  \begin{tabular}{@{}lll@{}}
    \toprule
    Parameter & Symbol & Value \\
    \midrule
    Max steer angle        & $\theta_{\max}$               & $0.5585$\,rad ($\approx 32^{\circ}$) \\
    Steering PD $K_p$      & $K_p$                          & $1600$\,N$\cdot$m/rad \\
    Steering PD $K_d$      & $K_d$                          & $160$\,N$\cdot$m$\cdot$s/rad \\
    Steering effort clamp  & $\tau_{\mathrm{steer,max}}$   & $1200$\,N$\cdot$m \\
    Drive torque           & $\tau_{\mathrm{drive,max}}$   & $619.8$\,N$\cdot$m \\
    Brake torque (front)   & $\tau_{\mathrm{brk,F}}$      & $1292.1$\,N$\cdot$m \\
    Brake torque (rear)    & $\tau_{\mathrm{brk,R}}$      & $638.5$\,N$\cdot$m \\
    Longitudinal scale     & $s_{\mathrm{long}}$           & $0.9273$ \\
    \bottomrule
  \end{tabular}
\end{table}


\subsection{Reward Design}
\label{sec:rewards}
\chattodo{Document the full reward signal: goal-reach bonus (sparse), route progress (dense), lane-keeping (geometric), collision penalty, road-edge TTC penalty, off-road penalty, idling penalty, TTC-based proximity penalty.  A reward component table with weight values would be ideal.  Mention that getting these weights right required extensive iteration — this is honest and useful.}
\chatcomment{Consider showing an ablation: what happens if you remove individual reward terms?  Even a small table showing ``full reward vs.\ no-TTC vs.\ no-route-progress'' would strengthen the paper.}


The per-step reward is the sum of ten terms: four sparse event rewards that also trigger agent termination, and six dense shaping rewards applied at every active step.
Table~\ref{tab:reward_terms} lists each component with its weight and key thresholds.
Tuning these weights required extensive iteration; the values reported here correspond to the configuration used in all experiments.

\paragraph{Sparse event rewards.}
Four mutually exclusive events terminate an agent and deliver a one-shot reward:
%
\begin{enumerate}[leftmargin=*, itemsep=1pt]
  \item \textbf{Goal reach} ($+45.0$).
  Awarded when the ego position is within $d_{\mathrm{goal}} = 3.0$\,m of the goal in XY.

  \item \textbf{Collision} ($-6.0$).
  Triggered when the maximum contact-sensor force exceeds 25\,N.
  A warm-up window of 24 steps after spawn suppresses false positives from initial settling.

  \item \textbf{Crash} ($-10.0$).
  Triggered when the vehicle falls below the ground plane ($z_{\mathrm{rel}} < 0$), drifts more than 100\,m from its start position, or tilts beyond a gravity-vector threshold ($g^b_z > -0.15$, indicating a rollover or severe pitch).

  \item \textbf{Lane-forbidden} ($-20.0$).
  Triggered when any of the vehicle's three bounding circles (the same proxy used for TTC, Section~\ref{sec:observations}) overlaps a road-edge segment (Waymo types 15--16), detected via an oriented-bounding-box overlap test.
\end{enumerate}
%
After termination, the agent is teleported off-stage and all subsequent dense rewards are masked to zero.

\paragraph{Dense shaping rewards.}
Six terms provide per-step signal for active agents:
%
\begin{enumerate}[leftmargin=*, itemsep=1pt]
  \item \textbf{Route progress} (weight $w = 2.0$).
  The displacement between consecutive steps is projected onto the tangent of the nearest driving-lane segment (Waymo types 1--2), oriented toward the goal:
  $r_{\mathrm{prog}} = \mathrm{clamp}\bigl((\mathbf{p}_t - \mathbf{p}_{t-1}) \cdot \hat{\mathbf{t}}_{\mathrm{lane}},\; -2,\; +2\bigr) \cdot w$.
  This rewards forward progress along the route and penalizes backwards motion.

  \item \textbf{Geometric lane-keeping} (weight $w = 0.08$).
  A continuous quality score combines lateral offset and heading alignment with the nearest driving lane:
  \[
    q = \exp\!\bigl(-(\ell / \sigma)^2\bigr) \cdot \bigl[(1 - w_h) + w_h \cdot \max(0,\, \cos(\psi - \psi_{\mathrm{lane}}))\bigr],
  \]
  where $\ell$ is the signed lateral offset, $\sigma = 1.75$\,m is the tolerance, and $w_h = 0.8$ is the heading weight.
  The reward is $r_{\mathrm{lane}} = q \cdot w$.

  \item \textbf{Off-road penalty} (weight $w = -0.5$).
  Applied when the vehicle's lateral offset exceeds 3.25\,m or its distance to the nearest lane exceeds 6.0\,m.

  \item \textbf{Idle penalty} (weight $w = -0.05$).
  Applied when the vehicle speed drops below 1.0\,m/s, discouraging the policy from learning to stop and wait.

  \item \textbf{Inter-vehicle TTC penalty.}
  The minimum pairwise TTC across all neighbors (using the swept-circle test from Section~\ref{sec:observations}) is converted to a penalty:
  \[
    p(\tau) = \min\!\bigl(\alpha_v / \max(\tau,\, \tau_{\mathrm{floor}}),\; p_{\mathrm{max},v}\bigr),
  \]
  with $\alpha_v = 0.10$, $p_{\mathrm{max},v} = 0.35$, and $\tau_{\mathrm{floor}} = 0.5$\,s.
  The reward is $r_{\mathrm{ttc}} = -p(\tau)$.

  \item \textbf{Road-edge TTC penalty.}
  The same $\alpha/\tau$ penalty form is applied to the closest road-edge segment within 40\,m ahead:
  $\alpha_e = 0.07$, $p_{\mathrm{max},e} = 0.40$, $\tau_{\mathrm{floor}} = 0.5$\,s.
  This encourages the policy to slow down when approaching road boundaries.
\end{enumerate}

\begin{table}[t]
  \centering\small
  \caption{Reward components used in all experiments.  Sparse rewards are delivered once and trigger agent termination; dense rewards are applied at every active step.}
  \label{tab:reward_terms}
  \begin{tabular}{@{}llrll@{}}
    \toprule
    Term & Type & Weight & Sign & Key threshold \\
    \midrule
    Goal reach         & Sparse & $45.0$  & $+$ & dist $\le 3.0$\,m \\
    Collision          & Sparse & $6.0$   & $-$ & force $\ge 25$\,N, 24-step warmup \\
    Crash              & Sparse & $10.0$  & $-$ & fall / drift $> 100$\,m / tilt \\
    Lane-forbidden     & Sparse & $20.0$  & $-$ & road-edge OBB overlap \\
    \midrule
    Route progress     & Dense  & $2.0$   & $\pm$ & clamp $[-2,+2]$\,m/step \\
    Lane-keeping       & Dense  & $0.08$  & $+$ & $\sigma{=}1.75$\,m, $w_h{=}0.8$ \\
    Off-road           & Dense  & $0.5$   & $-$ & lat $> 3.25$\,m or dist $> 6.0$\,m \\
    Idle               & Dense  & $0.05$  & $-$ & speed $< 1.0$\,m/s \\
    Inter-vehicle TTC  & Dense  & $\alpha{=}0.10$ & $-$ & max $= 0.35$, floor $= 0.5$\,s \\
    Road-edge TTC      & Dense  & $\alpha{=}0.07$ & $-$ & max $= 0.40$, floor $= 0.5$\,s \\
    \bottomrule
  \end{tabular}
\end{table}

\paragraph{Termination and time-out.}
An agent terminates upon any sparse event (goal, collision, crash, or lane-forbidden).
If no event occurs, the episode ends at $T_{\mathrm{max}} = 50$\,s (1{,}500 steps at 30\,Hz).
Terminated agents are teleported to an off-stage location ($x = 1{,}000$\,m) and excluded from subsequent reward and observation computation.

\chattodo{\textbf{ACTION ITEM --- ABLATION:} Consider a small ablation table showing goal-reach rate under the full reward vs.\ removing individual terms (e.g., no TTC penalty, no route progress).  Even three rows would strengthen the paper.}


% \subsection{Training Curriculum}
% \label{sec:curriculum}
% \chattodo{Describe the curriculum design used in training.  Document which reward terms are active in each stage and which checkpoints warm-start later stages.  Explain \emph{why} a curriculum is helpful: training with the full reward from scratch may fail because penalty signals dominate before the agent experiences any goal-reaching success.  The number and content of stages is still TBD; update once experiments are finalized.}
% \chatcite{Schulman et al., 2017 (PPO); Rudin et al., 2022 (RSL-RL / massively parallel RL).}

\subsection{Simulator Limitations and Sharp Edges}
%currenlty, if the scene has  roades crossing over each other, for example, overlapping road lane because of a bridge or ramp (i will include a image of that, ) the logic will break, as we compress the space into z=0. Secondly, as GPUdrive stated, in waymo dataset, the entrance of parking lot is labeled road edge. There exists cars that crossing the road edge to enter such area, so the Origin and destination points happen to appear on different sides of the road edge. These cases are going to introduce inconsistaenies and confusions in the training, as the collision with the road edge is unavoidable on the way , and the car is designed to get hard penalized for this behavior. In order sovle this, we coded a little tool to visualize and build the training scenes from a set of scenses, so the research can cherry pick out the scenes with such crossing. (specify this somewhere? ) This function is inclucded in a config wizard which helps you build a triaing configuration from scratch, accodring to the need. 
\label{sec:sharp_edges}
\chattodo{Document known issues: PhysX substep trade-offs, vehicle spawn filtering edge cases, any Isaac Lab quirks, scenarios where the physics breaks down, etc.  This section builds credibility.}


Transparency about simulator assumptions is essential for interpreting the results that follow.
Below we catalogue the most consequential simplifications, known failure modes, and practical workarounds in SceneFactory, organized by subsystem.

\paragraph{Z-flattening.}
The Waymo Open Motion Dataset provides full 3D polyline coordinates, including road elevation changes across bridges, ramps, overpasses, and multi-level interchanges.
SceneFactory collapses all road geometry to a constant ground height ($z{=}0$), because vehicles drive on a single flat PhysX ground plane per world.
This design decision has two consequences.
First, scenes with elevated structures (bridges, highway overpasses) produce visually overlapping road segments at ground level, which can confuse the $k$-nearest-road-point observation by presenting the agent with road points from an overhead or underpass road that it cannot physically reach.
Second, ramp grades and banked curves are lost, removing a class of driving challenges (hill starts, crest reactions) that affect real-world vehicle dynamics.
We mitigate the first issue through scene curation (see below) and by filtering road points whose original Waymo $z$ coordinate deviates significantly from the ground plane, but the fundamental limitation remains: SceneFactory is a 2.5D simulator operating on 3D data.

\paragraph{Scene curation.}
Not all Waymo scenarios are suitable for flat-ground training.
We developed an interactive \emph{Config Wizard} tool (\texttt{scene\_factory\_config\_wizard.py}) that presents each candidate scene as a bird's-eye plot showing lane-center polylines (Waymo types 1--2, blue), road-edge polylines (types 15--16, red), and valid start--goal pairs connected by origin--destination lines.
A human reviewer accepts or rejects each scene via keyboard input (y/n/b/q), with undo support.
The tool supports both a lightweight Matplotlib mode for rapid triage and an Isaac Sim preview mode that renders the full USD road geometry with 3D start/goal markers.

From a pool of ${\sim}100$ converted Waymo scenarios, we curated 64 training scenes and 64 held-out evaluation scenes, rejecting 23 training candidates and 49 evaluation candidates.
Common rejection reasons include: (i)~multi-level interchanges where z-flattening produces overlapping geometry that would poison the road-point observation, (ii)~scenes where origin--destination pairs straddle a road-edge polyline (see below), making the lane-forbidden penalty unavoidable, and (iii)~scenes with fewer than two valid controllable agents after spawn filtering.
The wizard produces a complete training configuration (scene list, environment count, reward weights, and curriculum stage) from scratch, so a researcher can go from raw Waymo data to a reproducible training run with minimal manual editing.
The curated 64-scene training set is tiled with random fill (seed 42) to populate all 256 parallel environments, so each scene appears in approximately four environments.

\paragraph{Waymo road-edge labeling.}
SceneFactory treats Waymo polyline types 15 and 16 as forbidden road edges for both the lane-forbidden termination event and the road-edge TTC penalty.
In practice, these labels are not uniformly reliable across the dataset.
As noted by Kazemkhani et al.~\cite{kazemkhani_gpudrive_2025}, the Waymo dataset labels parking-lot entrances as road edges.
Some tracked vehicles in these scenarios have their origin on the main road and their destination inside the parking area (or vice versa), so the recorded trajectory necessarily crosses a road-edge polyline.
When such a scenario is used for training, the agent faces an impossible conflict: reaching the goal requires crossing a boundary that triggers the lane-forbidden penalty ($-20.0$) and immediate termination.
These cases inject contradictory supervision into the training signal, degrading both goal-reach rate and lane-safety metrics.
Other scenarios lack road-edge annotations entirely where a physical curb or barrier clearly exists, creating blind spots in the lane-safety reward.
Scene curation (described above) removes the worst cases, but residual label noise remains and likely contributes to the modest lane-forbidden violation rates observed in evaluation.

\paragraph{Moving-only agent assumption.}
The spawn filter retains only Waymo agents whose start and goal positions are separated by more than 3.0\,m (the goal-reached threshold).
Agents below this threshold are treated as ``already at goal'' and excluded entirely.
Because the Waymo dataset records start and end positions from the logged trajectory, vehicles that remain parked, idle, or nearly stationary throughout the 9.1\,s recording window have near-zero start--goal displacement and are filtered out.
As a result, \emph{parked vehicles do not exist in the simulation}: a scene in which a moving vehicle must navigate past a row of curbside-parked cars will contain only the moving vehicle.
This simplification removes an entire class of realistic interactions (door-zone avoidance, double-parked obstacle negotiation, parallel parking maneuvers) and inflates the effective lane width available to the agent.
A future extension could retain filtered agents as static rigid-body obstacles rather than discarding them.

\paragraph{PointInstancer instability under Fabric.}
Road segments are rendered as USD \texttt{PointInstancer} prims for compactness (one instancer per road type per world, with per-instance position, orientation, and scale arrays).
We discovered that PointInstancer prims interact poorly with Isaac Lab's Fabric state management layer: visibility updates and prototype-table initialization are unreliable, producing intermittent rendering artifacts (invisible or misplaced road segments) that do not affect physics but corrupt visual debugging and video captures.
An explicit-prim fallback mode is provided (one \texttt{UsdGeom.Cube} per road segment) for diagnostic and stage-export workflows.
Production training disables Fabric entirely (\texttt{use\_fabric: false}) to avoid related issues with road metadata access, at a modest throughput cost.

\paragraph{Collision-sensor warmup.}
When vehicles are spawned or teleported to their start poses, PhysX requires several substeps to settle the suspension and establish stable wheel--ground contact.
During this transient, contact sensors report large spurious forces as the chassis drops onto the ground plane.
Without mitigation, these forces trigger false collision terminations at the start of every episode.
We suppress collision detection for the first 24 environment steps (${\approx}0.8$\,s) after each spawn, which is sufficient for the suspension to settle.
This warmup window introduces a brief period at episode start during which actual inter-vehicle collisions are unpenalized, though in practice vehicles are spawned with sufficient separation that early collisions are rare.

\paragraph{Teleport-only reset.}
Isaac Lab's default environment reset tears down and reconstructs PhysX actors, which is prohibitively expensive for multi-agent scenes with per-environment road geometry.
SceneFactory uses a \emph{teleport-only} reset strategy: after a single full initialization reset, subsequent episode resets reposition vehicles by writing directly to root-state tensors (position, orientation, linear and angular velocity) without rebuilding the physics scene.
This is ${\sim}10{\times}$ faster but has a subtle limitation: accumulated PhysX internal state (joint impulse caches, contact-pair broadphase data) is not fully cleared.
In rare cases this causes a vehicle to retain a residual velocity or contact state from the previous episode for one or two substeps after teleport.
The collision warmup window (above) masks most such artifacts.

\paragraph{Uniform friction surface.}
As discussed in Section~\ref{sec:weather_friction}, per-world friction is currently observation-only.
Even once per-vehicle wheel materials are implemented, friction will remain spatially uniform within each world: every wheel contact with the ground plane uses the same coefficient.
Real roads exhibit local friction variation (wet patches, oil spills, transitions between asphalt and concrete), and scenarios such as a slippery lane adjacent to a dry lane cannot be represented.
Modeling local friction would require either a segmented ground mesh with per-face materials or runtime material-property writes keyed to vehicle position, neither of which is straightforward under PhysX~5's current contact pipeline.

\paragraph{PhysX substep granularity.}
Physics runs at 120\,Hz with a decimation factor of 4, yielding 30\,Hz control.
The 120\,Hz rate is a compromise between simulation fidelity and throughput: higher rates (240\,Hz or 480\,Hz) would improve suspension stability and contact resolution but proportionally reduce training speed.
Our sysid calibration (Section~\ref{sec:vehicle_sysid}) was performed at this rate, so the fitted parameters absorb any 120\,Hz-specific artifacts; re-calibration would be needed if the substep rate were changed.

\paragraph{Clone-in-Fabric disabled.}
Isaac Lab's \texttt{clone\_in\_fabric} optimization, which clones environment state entirely in the Fabric layer for maximum GPU throughput, is disabled in SceneFactory because it is incompatible with the post-clone road injection workflow described in Section~\ref{sec:scene_construction}.
Fabric assumes that all environments are structurally identical after cloning; SceneFactory's per-environment road prims violate this assumption, causing metadata reads to return stale or incorrect values.
Disabling Fabric increases CPU--GPU synchronization overhead but is necessary for correctness.

\paragraph{Constant-velocity TTC.}
Both the inter-vehicle and road-edge TTC computations assume constant velocity over the prediction horizon.
This systematically overestimates collision risk during deceleration (the vehicle is predicted to continue at its current speed toward the obstacle) and underestimates risk during acceleration.
A quadratic (constant-acceleration) TTC formulation would be more accurate, particularly in approach zones where vehicles are actively braking, but would increase the per-step computational cost of the $K_v{\times}9$ pairwise circle tests.

% ============================================================
%  sections/experiments.tex  —  Experiments
% ============================================================
\section{Experiments}
\label{sec:experiments}

We evaluate SceneFactory along three axes:
simulation throughput (\S\ref{sec:throughput}),
end-to-end navigation performance on held-out Waymo scenes (\S\ref{sec:navigation_perf}),
and qualitative trajectory analysis (\S\ref{sec:qualitative}).
All experiments use a single shared navigation policy trained with PPO; the contribution of this work is the simulation platform itself, so we report policy performance as an end-to-end validation of the pipeline rather than as a claim of state-of-the-art driving.

% ---------------------------------------------------------------
\subsection{Experimental Setup}
\label{sec:exp_setup}
% ---------------------------------------------------------------

\paragraph{Hardware and software.}
All training and evaluation runs execute on a single NVIDIA RTX PRO 6000 Blackwell GPU (96\,GB VRAM).
The software stack consists of
NVIDIA Isaac Sim~\chattodo{version number},
Isaac Lab~\chattodo{version number}~\cite{nvidia_isaac_2025},
RSL-RL~\cite{schwarke_rsl-rl_2025} for PPO,
PyTorch~\chattodo{version}, and
CUDA~\chattodo{version}.
Physics runs at 120\,Hz (time step $\Delta t = 1/120$\,s);
the control policy acts at 30\,Hz (decimation factor 4).

\paragraph{Scene curation.}
From approximately 1{,}000 Waymo Open Motion Dataset scenarios~\cite{Kan_2024_icra},
we extract candidate road topologies through the pipeline described in \S\ref{sec:data_processing}.
A human operator first visually inspects each candidate scene with a curation wizard that renders lane-center polylines, road-edge polylines, and origin-destination pairs.
The operator accepts scenes that have plausible road geometry and rejects those with degenerate lane structure, multi-level overlaps from z-flattening, or unreachable goals (see \S\ref{sec:sharp_edges} for rejection criteria).
This produces \textbf{64 training scenes} and a disjoint set of \textbf{64 held-out test scenes}.
Within each accepted scene, an automated filter then retains only those agents whose start and goal positions both lie within the 400\,m $\times$ 400\,m world bounding box, are separated by at least 3.0\,m, and have valid coordinates.
Up to 16 qualifying agents are extracted per scene; scenes with fewer valid agents leave the remaining agent slots inactive.

\paragraph{Training configuration.}
Each training scene is replicated four times to fill $64 \times 4 = 256$ parallel environments, each hosting up to 16 agent slots, for a nominal capacity of 4{,}096 agents.
In practice, many Waymo scenes contain fewer than 16 qualifying agents after the spawn filter described below, so slots without a valid spawn remain inactive throughout the episode.
Worlds are assigned to environments by random shuffling (seed~42), so every PPO update samples transitions from all 64 unique road topologies.
Episodes last 50\,s (1{,}500 control steps).
Upon reaching the goal (within 3.0\,m) or triggering a termination condition (collision, off-road crash, timeout), an agent is parked in place with zero velocity; the remaining agents continue until every agent in the environment has terminated or the episode times out, at which point the entire world resets with fresh start-goal assignments for all agents.

\paragraph{PPO hyperparameters.}
Training uses PPO with an adaptive learning rate (initial $2.09{\times}10^{-4}$, KL target 0.01), rollout horizon of 128 steps, 5 gradient epochs over 16 mini-batches, $\gamma{=}0.99$, $\lambda{=}0.98$, and clip $\epsilon{=}0.15$.
The total batch size per update is $4{,}096 \times 128 = 524{,}288$ transitions.
The full hyperparameter table is in Appendix~\ref{app:ppo_hparams}.

\paragraph{Evaluation protocol.}
We evaluate on the 64 held-out test scenes, each loaded into a separate environment with up to 16 agent slots.
Across the 64 test worlds, 624 of the 1{,}024 nominal slots contain a valid spawn (the remainder correspond to scenes with fewer than 16 qualifying agents).
The policy runs in deterministic mode (mean of the action distribution).
Evaluation uses \emph{invincible mode}: collisions and off-road events are logged but do not terminate the agent, allowing it to continue toward the goal so that we can measure the unconditional goal-reach rate.
An agent is counted as successful if it reaches within 3.0\,m of its assigned goal before the 50\,s episode deadline.

% ---------------------------------------------------------------
\subsection{Throughput}
\label{sec:throughput}
% ---------------------------------------------------------------

A core engineering contribution of SceneFactory is throughput: training with full rigid-body dynamics at a rate that makes large-scale PPO practical on a single GPU.
To quantify this, we compare SceneFactory against its direct predecessor, a pipeline built on the same Isaac Sim simulator and PhysX 5 solver but using the built-in Vehicle Wizard vehicle model and Stable-Baselines3 for PPO (hereafter ``Baseline'').
Both pipelines load the same Waymo scene, use 16 agents per world, identical observation features (350 road points, 24 vehicle neighbors with TTC), and run on the same RTX PRO 6000 Blackwell GPU.
The only differences are the vehicle representation, the environment wrapper architecture, and the RL framework.

\paragraph{Metric.}
We report \emph{controlled agent steps per second} (CASPS): the number of policy-controlled agents that complete one 30\,Hz control tick per wall-clock second.
Both pipelines log CASPS every PPO iteration via embedded sub-phase timers;
we discard the first 50 logged iterations as warmup and report the mean and standard deviation over the remaining steady-state samples.

\paragraph{Scaling experiment.}
We select a single Waymo scene with 28 qualifying agents (ensuring all 16 slots are filled in every world) and replicate it at four scales: 32, 64, 128, and 256 parallel worlds, yielding 512 to 4{,}096 agent slots.
Each configuration trains for a short benchmark window: 500\,K total timesteps for the Baseline, 20 PPO iterations for SceneFactory.
Figure~\ref{fig:throughput_scaling} plots CASPS as a function of total agent slots for both pipelines.
The Baseline throughput is essentially flat at $\sim$152--174 CASPS across all world counts, confirming that per-agent Python loops---not the PhysX solver---are the bottleneck.
SceneFactory scales near-linearly, reaching 19{,}250 CASPS at 256 worlds (4{,}096 agent slots): a \textbf{127$\times$ speedup} over the Baseline at the same scale.
Even at the smallest configuration (32 worlds), SceneFactory is 22$\times$ faster.
Neither pipeline exceeded GPU memory at 256 worlds on the RTX PRO 6000 Blackwell (96\,GB VRAM).

\begin{table}[t]
\centering
\caption{Throughput scaling (CASPS, mean $\pm$ std over steady-state iterations) on a single RTX PRO 6000 Blackwell GPU.
Both pipelines use the same Waymo scene, the same PhysX 5 rigid-body solver, and 16 agents per world.}
\label{tab:throughput_scaling}
\small
\begin{tabular}{rrcc}
\toprule
Worlds & Agent slots & Baseline & SceneFactory \\
\midrule
 32 &   512 & $174 \pm 27$ & $3{,}870 \pm 66$ \\
 64 & 1{,}024 & $159 \pm 29$ & $7{,}173 \pm 114$ \\
128 & 2{,}048 & $164 \pm 1$ & $12{,}225 \pm 1{,}520$ \\
256 & 4{,}096 & $152 \pm 0$ & $19{,}250 \pm 2{,}984$ \\
\bottomrule
\end{tabular}
\end{table}

\paragraph{Why the speedup.}
The speedup arises from three architectural changes: (1)~GPU-resident joint-state tensors that replace per-agent Python loops with batched \texttt{torch} operations, (2)~cloned PhysX solver islands that enable parallel broad-phase/narrow-phase dispatch, and (3)~an on-device training loop (RSL-RL) that eliminates CPU--GPU data transfers.
A detailed per-phase timing breakdown (physics, observations, rewards, actions) is provided in Appendix~\ref{app:throughput}.

\paragraph{Comparison with kinematic simulators.}
Table~\ref{tab:cross_simulator} places SceneFactory in the context of published throughput numbers from kinematic driving simulators.
GPUDrive~\cite{kazemkhani_gpudrive_2025} (ICLR 2025) reports a peak of 2.3\,M agent-steps per second (ASPS) on an A100 GPU using a kinematic bicycle model.
When restricted to \emph{controlled} agents---i.e., those with non-trivial goals---the effective rate drops to $\sim$200\,K CASPS due to the high fraction of parked or near-goal vehicles in randomly sampled Waymo scenes (mean $\sim$10.8 controllable agents per scene out of up to 128 agent slots).
During end-to-end PPO training the authors report 200\,K--500\,K CASPS depending on hardware and observation configuration.
Nocturne~\cite{vinitsky_nocturne_2022} (NeurIPS 2022), a CPU-based C++ simulator with a kinematic model, achieves $\sim$15\,K ASPS on 16 CPU cores.
Waymax~\cite{gulino_waymax_2023} (NeurIPS 2024), a JAX-based simulator, supports GPU acceleration but encountered out-of-memory errors beyond 16 parallel environments in the GPUDrive authors' benchmarks, so we do not list a peak throughput.

SceneFactory operates at a lower absolute throughput than GPUDrive because each agent is a 10-DOF articulated rigid body simulated by the PhysX 5 solver at 120\,Hz (four physics sub-steps per 30\,Hz control tick).
The comparison is not apples-to-apples: the physics fidelity that makes SceneFactory useful for sim-to-real transfer---tire slip, suspension dynamics, inertial coupling---is precisely what makes it slower than a kinematic model.
Nevertheless, at 19.3\,K CASPS SceneFactory already exceeds Nocturne's throughput while providing full rigid-body dynamics, and it operates within an order of magnitude of GPUDrive's training-time CASPS despite the fundamentally heavier physics workload.
The relevant same-hardware, same-physics comparison is against the Baseline pipeline on the same RTX PRO 6000, where SceneFactory achieves a \textbf{127$\times$} speedup (Table~\ref{tab:throughput_scaling}).

\begin{table}[t]
\centering
\caption{Cross-simulator throughput comparison using published numbers.
All kinematic simulators use simplified vehicle dynamics (no rigid-body contact, no tire model); SceneFactory is the only entry with full PhysX 5 rigid-body physics.
ASPS = agent steps per second (all agents, including parked); CASPS = controlled agent steps per second (policy-active agents only).
GPUDrive and Nocturne numbers are taken from~\cite{kazemkhani_gpudrive_2025}; SceneFactory numbers are measured on our hardware.}
\label{tab:cross_simulator}
\small
\begin{tabular}{llllr}
\toprule
Simulator & Physics model & GPU & Metric & Throughput \\
\midrule
GPUDrive~\cite{kazemkhani_gpudrive_2025}
  & Kinematic bicycle & A100 80\,GB & ASPS & 2{,}300{,}000 \\
GPUDrive (training)
  & Kinematic bicycle & A100 80\,GB & CASPS & 200K--500K \\
Nocturne~\cite{vinitsky_nocturne_2022}
  & Kinematic & CPU 16 cores & ASPS & 15{,}000 \\
\midrule
\textbf{SceneFactory (ours)}
  & \textbf{10-DOF rigid body} & \textbf{RTX PRO 6000} & \textbf{CASPS} & \textbf{19{,}250} \\
\bottomrule
\end{tabular}
\end{table}

% ---------------------------------------------------------------
\subsection{Navigation Performance}
\label{sec:navigation_perf}
% ---------------------------------------------------------------

We train a single shared navigation policy across all 4{,}096 agent slots (256 worlds $\times$ 16 agents) using PPO and select the best checkpoint by validation goal-reach rate.
The selected checkpoint (model\_975) corresponds to approximately $975 \times 128 \times 4{,}096 \approx 511\text{M}$ environment steps.

\paragraph{Evaluation protocol.}
All evaluations use \emph{invincible mode}: collision, lane-boundary, and off-road events are logged but do not terminate the agent, so that every spawned vehicle is observed until it either reaches its goal (within 3.0\,m) or the 50\,s episode budget expires.
This protocol allows us to measure goal-reach rates without survivorship bias.

\paragraph{Generalization study.}
To assess how well the learned policy generalizes across scenes and routing targets, we evaluate the same checkpoint under three conditions that progressively decouple the test setting from the training distribution:
%
\begin{enumerate}[leftmargin=*, itemsep=1pt]
  \item \textbf{Test scenes, original OD.}
        64 held-out Waymo scenes with the original origin--destination pairs extracted from the dataset.
  \item \textbf{Test scenes, random OD.}
        The same 64 held-out scenes, but each agent's destination is re-sampled uniformly along reachable road segments at 20--60\,m travel distance.
  \item \textbf{Train scenes, random OD.}
        The 64 unique training scenes, each with 16 randomly sampled origin--destination pairs (20--60\,m travel distance), for a total of 1{,}024 agents---matching the test-scene agent count.
\end{enumerate}

Table~\ref{tab:generalization} reports the results.

\begin{table}[t]
\centering
\caption{Generalization evaluation of the navigation policy (model\_975).
All runs use invincible mode with a 50\,s episode budget and 3.0\,m goal threshold.
``Original OD'' denotes destination points extracted from the Waymo dataset;
``Random OD'' denotes destinations re-sampled uniformly at 20--60\,m travel distance along reachable road segments.}
\label{tab:generalization}
\small
\begin{tabular}{llccccc}
\toprule
Scenes & OD & Agents & Success $\uparrow$ & Collision $\downarrow$ & Lane-Forbidden $\downarrow$ & $\bar{d}_\text{goal}$ (m) $\downarrow$ \\
\midrule
Test (64)   & Original & 624   & 71.6\% & 6.3\% & 9.0\% & 10.1 \\
Test (64)   & Random   & 1{,}024 & 92.8\% & 4.5\% & 6.9\% & 3.7 \\
Train (64)  & Random   & 1{,}024 & 92.5\% & 3.5\% & 7.6\% & 3.7 \\
\bottomrule
\end{tabular}
\end{table}

On held-out test scenes with original destinations, the policy reaches its goal in 71.6\% of cases (447 of 624 agents) with a mean residual distance of 10.1\,m for non-successful agents.
When the same test scenes are paired with randomly sampled destinations (Condition~2), the success rate jumps to 92.8\% (950 of 1{,}024 agents), nearly matching the 92.5\% achieved on training scenes with the same random-OD protocol (Condition~3).
The near-identical performance between Conditions~2 and~3 is the central finding: \emph{the policy generalizes to entirely novel road geometries as well as it performs on familiar ones}, when the routing difficulty is controlled.
This rules out scene memorization and confirms that the policy has learned a general goal-conditioned navigation strategy.

The 21-percentage-point gap between Conditions~1 and~2 (71.6\% vs.\ 92.8\%), which share the same test scenes but differ only in OD pairs, reveals that the original Waymo destinations are substantially harder than the 20--60\,m random-OD protocol.
The Waymo-extracted routes include long-range multi-turn trajectories, goals near road boundaries, and occasionally ambiguous lane assignments that the random sampler avoids by construction.
This suggests that further gains on the original OD setting are achievable through longer episodes, curriculum over route difficulty, or improved goal-approach behavior, rather than additional scene diversity.

The collision rates remain low across all conditions (3.5--6.3\%), and are inflated by invincible mode, which allows post-collision agents to continue driving and potentially trigger additional collision events; under normal termination semantics each collision would end the episode immediately.
Notably, the off-road crash rate is higher on test scenes (18.7\%) than on training scenes (14.8\%) under the same random-OD protocol, indicating that novel road geometries do cause more lane departures---but the policy recovers and still reaches the goal at nearly the same rate.

\paragraph{Comparison with kinematic simulators.}
GPUDrive reports near-perfect goal-reach rates ($>$98\%) on Waymo scenes using a kinematic bicycle model with no rigid-body dynamics.
We attribute the gap primarily to the harder control problem posed by full PhysX dynamics: the policy must learn to manage tire slip, suspension oscillations, and inertial overshoot that do not exist in kinematic models.
The 10-DOF articulated vehicle with effort-based actuation (\S\ref{sec:action_space}) requires the agent to implicitly discover braking distances, cornering limits, and steering lag, whereas a kinematic model executes velocity commands with no delay or slip.
Closing this gap through longer training, refined reward shaping, and curriculum design is active work; the current results demonstrate that the platform successfully trains navigating agents end-to-end under physically realistic dynamics.

\chattodo{Add a training curve figure (mean reward and/or success rate vs.\ training iteration) for the best run.  Even a single curve is informative.  Export from TensorBoard.}

% ---------------------------------------------------------------
\subsection{Qualitative Analysis}
\label{sec:qualitative}
% ---------------------------------------------------------------

\chattodo{Include 3--4 bird's-eye trajectory plots from the evaluation.  The eval run already produced \texttt{videos/scene\_factory\_policy\_eval.mp4} and per-step JSONL logs that can be rendered.  Show: (1) a world where all agents succeed with clean lane-following, (2) a dense world with successful multi-agent coordination, (3) a failure case (e.g., the worst-performing world with all collisions).  Annotate start positions (green), goals (yellow), and agent trajectories (colored lines).}

Figure~\ref{fig:qualitative_trajectories} shows bird's-eye trajectory visualizations from three representative test worlds.
\chattodo{Create and include the figure.}

\paragraph{Success modes.}
In the majority of test worlds, agents follow lane centers, decelerate smoothly near goals, and avoid collisions with neighboring vehicles.
The policy learns to yield at merge points when a neighbor vehicle occupies the target lane, demonstrating emergent multi-agent coordination despite using only local observations (24 nearest neighbors within the observation radius).

\paragraph{Failure modes.}
We identify three recurring failure patterns:
\begin{enumerate}[leftmargin=*, itemsep=2pt]
  \item \textbf{Overshoot at sharp turns.}
        In worlds with tight intersections or U-turns, agents carry too much speed into the turn, overshoot the lane boundary, and trigger an off-road crash.
        This failure is absent in kinematic simulators where instantaneous heading changes are possible.
  \item \textbf{Goal approach oscillation.}
        Some agents oscillate near the goal without entering the 3.0\,m acceptance radius.
        The sign-memory brake mechanism (\S\ref{sec:action_space}) mitigates this compared to earlier training runs, but a small fraction of agents still fail to decelerate early enough.
  \item \textbf{Dense-scene congestion.}
        In worlds with all 16 agent slots occupied, vehicles occasionally deadlock when two agents attempt to merge into the same lane segment simultaneously.
        The constant-velocity TTC proxy (\S\ref{sec:sharp_edges}) provides limited lookahead for such scenarios.
\end{enumerate}

\chattodo{Verify these failure modes against the per-step JSONL logs from the eval run.  Quantify: what fraction of failures falls into each category?}

% ============================================================
%  sections/conclusion.tex  —  Conclusion and Future Work
% ============================================================
\section{Conclusion}
\label{sec:conclusion}

% ---------------------------------------------------------------
% ORIGINAL SOC TEXT — preserved verbatim
% ---------------------------------------------------------------

A conclusion, which should be good when we have eveyrthing ready.

% ---------------------------------------------------------------
%  Advisor annotations
% ---------------------------------------------------------------

\chatstructure{The conclusion should have two parts: (1) a summary of what was accomplished, and (2) future work.  Draft outline below.}

\begin{chatbox}[Suggested Conclusion Outline]
\textbf{Summary paragraph:}
\begin{itemize}
  \item We presented SceneFactory, a GPU-accelerated, physics-enabled, multi-agent driving simulator built on NVIDIA Isaac Sim and Isaac Lab.
  \item SceneFactory procedurally constructs multi-world training environments from Waymo Open Motion Dataset scenarios, with full PhysX 5 rigid-body dynamics and a weather-to-friction module.
  \item We demonstrated that the platform achieves [X] env-steps/sec with [Y] parallel vehicles on a single GPU, and that the three-stage curriculum enables end-to-end navigation policy learning.
  \item Restate the gap filled: first system combining GPU-parallel multi-world execution, multi-agent support, full physics, and weather-aware friction.
\end{itemize}

\textbf{Limitations paragraph:}
\begin{itemize}
  \item Current navigation success rates do not match kinematic-dynamics simulators; we attribute this to the harder control problem posed by full rigid-body dynamics and plan to address it through [specific plans: longer training, improved reward shaping, curriculum refinements].
  \item Scene diversity is currently limited to a curated Waymo subset; scaling to the full dataset is engineering work in progress.
\end{itemize}

\textbf{Future work paragraph (keep brief --- 3--4 sentences max):}
\begin{itemize}
  \item Adaptive robotic work zone management as a co-trained agent.
  \item Integration of formal safety metrics (TTC distributions, PET, DRAC).
  \item Dynamic, spatially varying weather conditions within episodes.
  \item SceneFactory-to-Cosmos pipeline for photorealistic video generation.
\end{itemize}
\end{chatbox}

\chattodo{Write the actual conclusion once experimental results are finalized.  Keep it to $\sim$0.5 pages.}


% ============================================================
%  Acknowledgments (hidden during review; shown in camera-ready)
% ============================================================
\begin{ack}
  % Fill in after acceptance.
  % e.g., This work was supported by …
\end{ack}

% ============================================================
%  References
% ============================================================
\bibliographystyle{plainnat}
\bibliography{references}

% ============================================================
%  NeurIPS Paper Checklist  (required — do NOT remove)
% ============================================================
\input{checklist}

% ============================================================
%  Appendix
% ============================================================
% ============================================================
%  sections/appendix.tex  —  Supplementary Material (Appendix)
% ============================================================
\appendix

\section{Weather-to-Friction Model Details}
\label{app:friction_model}

This appendix provides the full mathematical formulation of the tire--pavement friction model used in SceneFactory (summarized in \S\ref{sec:weather_friction}).
We implement the modified Average Lumped LuGre (ALL) model of Zhao et al.~\cite{zhao_tire-pavement_2024}, which extends the LuGre bristle-deflection framework to account for pavement texture and water film thickness.

\paragraph{Bristle dynamics.}
The model evolves the average bristle deflection $\bar{z}$ via the ODE
\begin{equation}
\label{eq:app_bristle_ode}
\frac{d\bar{z}}{dt} = v_r - \theta \, Y_R \left[\frac{\sigma_0 \lvert v_r \rvert}{g(v_r)}\,\bar{z} + K \lvert \omega_r \rvert\,\bar{z}\right],
\end{equation}
where $v_r$ is the relative (slip) speed between the tire contact patch and the road, $\omega_r$ is the wheel circumferential speed, $\sigma_0$ is the bristle stiffness, $K = 7/(6L)$ is the ALL approximation constant with $L$ the contact patch length, $\theta$ is a texture influence coefficient (per road surface type), and $Y_R \in [0,1]$ is the contact patch length ratio under hydrodynamic lift.

\paragraph{Stribeck function.}
The steady-state function $g(v_r) = \mu_c + (\mu_s - \mu_c)\exp\bigl(-\lvert v_r / v_s \rvert^\alpha\bigr)$ captures the transition from static friction $\mu_s$ to Coulomb friction $\mu_c$.
The Stribeck speed depends on the water film thickness $h_w$ (in meters) via
\begin{equation}
\label{eq:app_stribeck_speed}
v_s(h_w) = b_1 \exp(1000\,b_2\,h_w + b_3) + b_4,
\end{equation}
where $b_1 = 4.8916$, $b_2 = -7.91$, $b_3 = 3.01$, and $b_4 = 3.40$ are calibrated constants from~\cite{zhao_tire-pavement_2024}.
Because $b_2 < 0$, $v_s$ decreases with water film thickness: thicker water films cause the friction drop-off to onset at lower slip speeds.

\paragraph{Effective friction coefficient.}
Given the steady-state solution $\bar{z}^*$, the effective friction coefficient is
\begin{equation}
\label{eq:app_mu_all}
\mu = \max\!\bigl(\theta\,Y_R - Y_F,\; 0\bigr) \cdot \bigl(\theta\,Y_R\,\sigma_0\,\bar{z}^* + \sigma_1\,\dot{\bar{z}} + \sigma_2\,v_r\bigr),
\end{equation}
where $Y_F \in [0,1]$ is the hydrodynamic force ratio and $\sigma_1$, $\sigma_2$ are viscous damping coefficients.
Under the default calibration from~\cite{zhao_tire-pavement_2024}, both $\sigma_1 = 0$ and $\sigma_2 = 0$, so the effective formula reduces to $\mu = \max(\theta\,Y_R - Y_F,\,0)\cdot\theta\,Y_R\,\sigma_0\,\bar{z}^*$.
Both $Y_R$ and $Y_F$ depend on vehicle speed, water film thickness, tire geometry (contact patch dimensions, tread radius, wheel radius), and physical constants (water density and viscosity).
When the contact factor $\theta\,Y_R - Y_F$ drops to zero, the model predicts full hydroplaning.
Our implementation numerically integrates Equation~\ref{eq:app_bristle_ode} to steady state using exponential integration ($\Delta t = 10^{-3}$\,s, convergence threshold $\epsilon = 10^{-7}$), matching the paper's calibrated parameters (Table~4 and Table~5 in~\cite{zhao_tire-pavement_2024}).

\paragraph{Road surface presets.}
The texture influence coefficient $\theta$ and the texture amplitude $T_d$ vary by pavement type.
We provide three presets calibrated to the values in~\cite{zhao_tire-pavement_2024}:
\emph{Asphalt Concrete} (AC, $\theta{=}1.00$, $T_d{=}0.65$\,mm),
\emph{Stone Mastic Asphalt} (SMA, $\theta{=}1.09$, $T_d{=}0.80$\,mm), and
\emph{Open-Graded Friction Course} (OGFC, $\theta{=}1.21$, $T_d{=}1.08$\,mm).
OGFC, designed for drainage, retains higher friction under wet conditions; AC, the most common pavement, serves as the baseline.

\paragraph{PhysX integration pipeline.}
Each world in the training stage is assigned a friction configuration specifying the road surface type and the water film thickness $h_w$ (in mm).
At environment construction time, the pipeline feeds these values into the modified ALL model (Equation~\ref{eq:app_mu_all}) at a reference speed of 13.89\,m/s ($\approx$50\,km/h) and two slip ratios (0.15 for static, 0.80 for dynamic) to produce per-world static and dynamic friction coefficients $\mu_s$ and $\mu_d$.
The effective coefficient $\mu_{\mathrm{eff}}$ (static by default) is used downstream as follows.

The ground-plane material uses the sysid-tuned friction scale for the dry-asphalt surface type (Section~\ref{sec:vehicle_sysid}), with static friction $\mu^{\mathrm{ground}}_s = \min(1.0,\; f_{\mathrm{surface}} \cdot \sqrt{f_{\mathrm{lon}} \cdot f_{\mathrm{lat}}})$ and dynamic friction $\mu^{\mathrm{ground}}_d = 0.95\,\mu^{\mathrm{ground}}_s$.
PhysX resolves contact friction using the minimum of the ground and tire material values (combine mode \texttt{min}).
Because all cloned environments share a single global ground plane, this material is uniform across worlds.
Per-world friction variation is currently communicated to the policy only through the observation (Equation~\ref{eq:weather_obs}).

\chattodo{\textbf{ACTION ITEM --- FIGURE:} Consider a plot of $\mu_{\mathrm{eff}}$ vs.\ $h_w$ (0--1.0\,mm) for the three surface presets (AC, SMA, OGFC) at $v{=}$13.89\,m/s.}


% ============================================================
\section{Observation Space Details}
\label{app:observations}
% ============================================================

Each agent receives a 1{,}929-dimensional observation vector assembled from four groups.
All spatial quantities are expressed in the agent's body frame via a 2D rotation by $-\psi$ (ego heading), so the policy sees an ego-centric view of the world.

\paragraph{Ego state (7-dim).}
The ego state encodes the agent's goal and velocity in body-frame coordinates:
\begin{enumerate}[leftmargin=*, itemsep=1pt]
  \item Goal position $(g_x^b, g_y^b)$, each divided by the scene bounding-box half-size ($B/2 = 100$\,m).
  \item Heading error to the goal, encoded as $(\sin\Delta\psi,\,\cos\Delta\psi)$.
  \item Euclidean goal distance $\|g^b_{xy}\|$, divided by $B/2$.
  \item Body-frame velocity $(v_x^b, v_y^b)$, each divided by 10\,m/s.
\end{enumerate}
When the weather-to-friction module is active, the 4-dim weather context (Equation~\ref{eq:weather_obs}) is appended to the ego group, yielding an 11-dim input to the ego encoder.

\paragraph{Road geometry context ($K_r \times 5 = 1{,}750$-dim).}
Each agent selects $K_r = 350$ road-segment midpoints from the pre-loaded GPU metadata tensors (\S\ref{sec:data_processing}).
The selection operates in two steps: first, all road points within a radius $r = 10$\,m of the ego position are identified via a squared-distance threshold; second, these candidate points are sorted by their original polyline index to preserve road topology, and the first $K_r$ are taken.
If fewer than $K_r$ points fall within the radius, the remaining slots are zero-padded.
Each selected point contributes five features:
\begin{enumerate}[leftmargin=*, itemsep=1pt]
  \item Relative position $(\Delta x^b, \Delta y^b)$ in ego body frame, each divided by $r$.
  \item Polyline type code, divided by a normalizer ($= 20$).
  \item Road direction unit vector $(d_x^b, d_y^b)$ in ego body frame.
\end{enumerate}

\paragraph{Neighbor vehicles ($K_v \times 7 = 168$-dim).}
The $K_v = 24$ nearest alive agents (by Euclidean distance in world XY) are selected per ego agent.
Self-distance and dead-agent distances are set to infinity before the argsort.
Each neighbor contributes seven features:
\begin{enumerate}[leftmargin=*, itemsep=1pt]
  \item Relative position $(\Delta x^b, \Delta y^b)$ in ego body frame, each divided by $B/2$.
  \item Chassis length and width, each divided by $B/2$.
  \item Relative heading $\Delta\psi$, divided by $\pi$.
  \item Speed $\|v_{xy}\|$, divided by 10\,m/s.
  \item Pairwise time-to-collision (TTC), computed via a swept-circle test.
    Each vehicle is approximated by three equal-radius circles placed along the longitudinal axis at offsets $\{-d,\, 0,\, +d\}$ from the chassis centre, where $r = \max(0.45,\; 0.55\, W)$ and $d = \min\!\bigl(\tfrac{\mathit{WB}}{2},\; \max(0,\; \tfrac{L}{2} - 0.8\,r)\bigr)$.
    For the default vehicle ($L{=}4.0$\,m, $W{=}2.0$\,m, $\mathit{WB}{=}2.6$\,m) this gives $r \approx 1.10$\,m and $d \approx 1.12$\,m.
    The test solves a quadratic closest-approach equation for every pair among the $3 \times 3 = 9$ circle combinations and takes the minimum.
    The resulting TTC is clamped to $[0,\, \tau_{\max}]$ and divided by $\tau_{\max} = 10$\,s.
    Neighbors behind the ego or with no collision trajectory receive $\mathrm{TTC}/\tau_{\max} = 1$.
\end{enumerate}
Slots beyond the number of alive neighbors are zero-padded.


% ============================================================
\section{Network Architecture Details}
\label{app:network}
% ============================================================

The observation groups are processed by a late-fusion actor-critic network with per-modality encoders and max-pool set aggregation, following the design of GPUDrive~\cite{kazemkhani_gpudrive_2025}.
We make two modifications: (i)~we use \emph{masked} max-pool (zero-padded slots are filled with $-\infty$ before the $\max$) rather than a plain pool, and (ii)~we replace their Tanh activations with ELU.
The flat observation is split into its three groups and processed by separate encoder MLPs, one set for the actor and one for the critic (no weight sharing).

\begin{itemize}[leftmargin=*, itemsep=1pt]
  \item \textbf{Ego encoder}: $11 \to 64 \to 64$ (two hidden layers, ELU activations).
  \item \textbf{Road encoder}: $5 \to 96 \to 96$, applied independently to each of the $K_r$ points; the $K_r$ output vectors are aggregated via \emph{masked max-pool}, where zero-padded points are filled with $-\infty$ before the $\max$ operation.
  \item \textbf{Vehicle encoder}: $7 \to 96 \to 96$, applied independently to each of the $K_v$ neighbors; aggregated via masked max-pool in the same way.
\end{itemize}

The three embeddings (64 + 96 + 96 = 256) are concatenated and passed through a shared trunk MLP ($256 \to 128 \to 64$, ELU).
The actor head projects to a 3-dim mean vector $\boldsymbol{\mu}$ (one per action), paired with a state-independent learnable log-standard-deviation $\log\boldsymbol{\sigma}$ to form a diagonal Gaussian policy.
The critic head projects to a scalar state-value estimate.
No empirical observation normalization (running mean/std) is used; all features are manually scaled at assembly time as described in Appendix~\ref{app:observations}.

\chattodo{\textbf{ACTION ITEM --- FIGURE:} Create a network architecture diagram showing the three encoder branches, max-pool aggregation, fusion trunk, and actor/critic heads.}


% ============================================================
\section{Action Space Details}
\label{app:actions}
% ============================================================

Each agent outputs a three-dimensional continuous action $\mathbf{a} = (a_{\mathrm{thr}},\, a_{\mathrm{str}},\, a_{\mathrm{brk}})$.
The policy produces values in $[-1,1]^3$; throttle and brake are rectified to $[0,1]$ before decoding, so a zero-mean Gaussian policy naturally idles.
All three channels are decoded into PhysX \emph{effort targets}; the steering PD loop described in Eq.\ref{eq:steer_pd} is computed in software at every physics substep.

\paragraph{Control frequency.}
Physics runs at $1/\Delta t_{\mathrm{phys}} = 120$\,Hz.
Actions are held constant over a decimation window of 4 substeps, yielding a control frequency of 30\,Hz.
Within each substep the steering PD reads the current joint state, so the effective PD loop rate is 120\,Hz.

\paragraph{Brake sign memory.}
Brake torque must oppose wheel rotation, so its sign is $-\mathrm{sign}(\dot{q}_{\mathrm{wheel}})$.
When $|\dot{q}_{\mathrm{wheel}}| < 10^{-4}$\,rad/s the sign becomes unreliable; we therefore latch the last observed sign and continue applying it until rotation resumes.
This prevents torque chatter near zero velocity.

\paragraph{Actuation parameters.}
Table~\ref{tab:app_action_params} lists the sysid-tuned values used in all experiments.
Drive torque is applied to the two front wheels only (FWD); brake torque is applied to all four wheels with separate front/rear magnitudes.
The sysid procedure additionally fits per-surface friction scales (longitudinal and lateral) that modulate the effective contact friction for each road surface type; these are relevant only for multi-surface scenarios (\S\ref{sec:weather_friction}).

\begin{table}[h]
  \centering\small
  \caption{Action-decoding parameters (sysid-tuned values).}
  \label{tab:app_action_params}
  \begin{tabular}{@{}lll@{}}
    \toprule
    Parameter & Symbol & Value \\
    \midrule
    Max steer angle        & $\theta_{\max}$               & $0.450$\,rad ($\approx 25.8^{\circ}$) \\
    Steering PD $K_p$      & $K_p$                          & $1839.5$\,N$\cdot$m/rad \\
    Steering PD $K_d$      & $K_d$                          & $110.5$\,N$\cdot$m$\cdot$s/rad \\
    Steering effort clamp  & $\tau_{\mathrm{steer,max}}$   & $1200$\,N$\cdot$m \\
    Drive torque           & $\tau_{\mathrm{drive,max}}$   & $600.7$\,N$\cdot$m \\
    Brake torque (front)   & $\tau_{\mathrm{brk,F}}$      & $1090.5$\,N$\cdot$m \\
    Brake torque (rear)    & $\tau_{\mathrm{brk,R}}$      & $980.7$\,N$\cdot$m \\
    Wheel mass             & $m_{\mathrm{wheel}}$          & $37.5$\,kg \\
    Wheel inertia scale    & $s_{\mathrm{inertia}}$        & $1.094$ \\
    Suspension stiffness   & $k_{\mathrm{susp}}$           & $1080.8$\,N/m \\
    Suspension damping      & $c_{\mathrm{susp}}$           & $2764.3$\,N$\cdot$s/m \\
    Yaw stability damping  & $\lambda_{\mathrm{yaw}}$      & $10.6$\,N$\cdot$m$\cdot$s/rad \\
    \bottomrule
  \end{tabular}
\end{table}


% ============================================================
\section{Reward Design Details}
\label{app:rewards}
% ============================================================

The per-step reward is the sum of ten terms: four sparse event rewards that also trigger agent termination, and six dense shaping rewards applied at every active step.
Tuning these weights required extensive iteration; the values reported here correspond to the configuration used in all experiments.

\paragraph{Sparse event rewards.}
Four mutually exclusive events terminate an agent and deliver a one-shot reward:
\begin{enumerate}[leftmargin=*, itemsep=1pt]
  \item \textbf{Goal reach} ($+45.0$).
  Awarded when the ego position is within $d_{\mathrm{goal}} = 3.0$\,m of the goal in XY.

  \item \textbf{Collision} ($-6.0$).
  Triggered when the maximum contact-sensor force exceeds 25\,N.
  A warm-up window of 24 steps after spawn suppresses false positives from initial settling.

  \item \textbf{Crash} ($-10.0$).
  Triggered when the vehicle falls below the ground plane ($z_{\mathrm{rel}} < 0$), drifts more than 100\,m from its start position, or tilts beyond a gravity-vector threshold ($g^b_z > -0.15$, indicating a rollover or severe pitch).

  \item \textbf{Lane-forbidden} ($-20.0$).
  Triggered when any of the vehicle's three bounding circles overlaps a road-edge segment (Waymo types 15--16), detected via an oriented-bounding-box overlap test.
\end{enumerate}
After termination, the agent is teleported off-stage and all subsequent dense rewards are masked to zero.

\paragraph{Dense shaping rewards.}
Six terms provide per-step signal for active agents:
\begin{enumerate}[leftmargin=*, itemsep=1pt]
  \item \textbf{Route progress} (weight $w = 2.0$).
  The displacement between consecutive steps is projected onto the tangent of the nearest driving-lane segment (Waymo types 1--2), oriented toward the goal:
  $r_{\mathrm{prog}} = \mathrm{clamp}\bigl((\mathbf{p}_t - \mathbf{p}_{t-1}) \cdot \hat{\mathbf{t}}_{\mathrm{lane}},\; -2,\; +2\bigr) \cdot w$.
  This rewards forward progress along the route and penalizes backwards motion.

  \item \textbf{Geometric lane-keeping} (weight $w = 0.08$).
  A continuous quality score combines lateral offset and heading alignment with the nearest driving lane:
  \[
    q = \exp\!\bigl(-(\ell / \sigma)^2\bigr) \cdot \bigl[(1 - w_h) + w_h \cdot \max(0,\, \cos(\psi - \psi_{\mathrm{lane}}))\bigr],
  \]
  where $\ell$ is the signed lateral offset, $\sigma = 1.75$\,m is the tolerance, and $w_h = 0.8$ is the heading weight.
  The reward is $r_{\mathrm{lane}} = q \cdot w$.

  \item \textbf{Off-road penalty} (weight $w = -0.5$).
  Applied when the vehicle's lateral offset exceeds 3.25\,m or its distance to the nearest lane exceeds 6.0\,m.

  \item \textbf{Idle penalty} (weight $w = -0.05$).
  Applied when the vehicle speed drops below 1.0\,m/s, discouraging the policy from learning to stop and wait.

  \item \textbf{Inter-vehicle TTC penalty.}
  The minimum pairwise TTC across all neighbors is converted to a penalty:
  \[
    p(\tau) = \min\!\bigl(\alpha_v / \max(\tau,\, \tau_{\mathrm{floor}}),\; p_{\mathrm{max},v}\bigr),
  \]
  with $\alpha_v = 0.10$, $p_{\mathrm{max},v} = 0.35$, and $\tau_{\mathrm{floor}} = 0.5$\,s.
  The reward is $r_{\mathrm{ttc}} = -p(\tau)$.

  \item \textbf{Road-edge TTC penalty.}
  The same $\alpha/\tau$ penalty form is applied to the closest road-edge segment within 40\,m ahead:
  $\alpha_e = 0.07$, $p_{\mathrm{max},e} = 0.40$, $\tau_{\mathrm{floor}} = 0.5$\,s.
  This encourages the policy to slow down when approaching road boundaries.
\end{enumerate}

\begin{table}[t]
  \centering\small
  \caption{Reward components used in all experiments.  Sparse rewards are delivered once and trigger agent termination; dense rewards are applied at every active step.}
  \label{tab:app_reward_terms}
  \begin{tabular}{@{}llrll@{}}
    \toprule
    Term & Type & Weight & Sign & Key threshold \\
    \midrule
    Goal reach         & Sparse & $45.0$  & $+$ & dist $\le 3.0$\,m \\
    Collision          & Sparse & $6.0$   & $-$ & force $\ge 25$\,N, 24-step warmup \\
    Crash              & Sparse & $10.0$  & $-$ & fall / drift $> 100$\,m / tilt \\
    Lane-forbidden     & Sparse & $20.0$  & $-$ & road-edge OBB overlap \\
    \midrule
    Route progress     & Dense  & $2.0$   & $\pm$ & clamp $[-2,+2]$\,m/step \\
    Lane-keeping       & Dense  & $0.08$  & $+$ & $\sigma{=}1.75$\,m, $w_h{=}0.8$ \\
    Off-road           & Dense  & $0.5$   & $-$ & lat $> 3.25$\,m or dist $> 6.0$\,m \\
    Idle               & Dense  & $0.05$  & $-$ & speed $< 1.0$\,m/s \\
    Inter-vehicle TTC  & Dense  & $\alpha{=}0.10$ & $-$ & max $= 0.35$, floor $= 0.5$\,s \\
    Road-edge TTC      & Dense  & $\alpha{=}0.07$ & $-$ & max $= 0.40$, floor $= 0.5$\,s \\
    \bottomrule
  \end{tabular}
\end{table}

\paragraph{Termination and time-out.}
An agent terminates upon any sparse event (goal, collision, crash, or lane-forbidden).
If no event occurs, the episode ends at $T_{\mathrm{max}} = 50$\,s (1{,}500 steps at 30\,Hz).
Terminated agents are teleported to an off-stage location ($x = 1{,}000$\,m) and excluded from subsequent reward and observation computation.


% ============================================================
\section{System Identification Details}
\label{app:sysid}
% ============================================================

The vehicle model has 20 tunable dynamics parameters: drive, brake (front/rear), and steering torques; steering PD gains and limits; suspension stiffness and damping; wheel mass and inertia scale; chassis center-of-mass offset; lateral and yaw damping coefficients; and per-surface friction scales for three surface types (dry, wet, gravel).
We fit these parameters using the cross-entropy method (CEM).

\paragraph{Reference data.}
A PhysX ``teacher'' vehicle (the Isaac Sim built-in PhysX Vehicle model, which includes a validated tire model, drivetrain, and suspension) executes 139 scripted maneuvers on a multi-surface patch track.
The maneuvers span five tiers: 17 longitudinal (throttle sweep 10--100\%, brake sweep, linear ramps), 64 lateral (step-steer and constant-radius at four throttle levels $\times$ four--five steer amplitudes $\times$ left/right), 18 combined (trail-brake at three throttle$\times$steer$\times$direction), 39 frequency-response (sinusoidal steering at multiple throttle/amplitude/frequency settings, plus chirp sweeps), and 1 surface transition across dry asphalt ($\mu_s{=}1.0$), wet asphalt ($\mu_s{=}0.75$), and gravel ($\mu_s{=}0.60$).
The comprehensive suite covers the full operating envelope from 10\% to 100\% throttle, ensuring that the fitted parameters generalize across the vehicle's speed and lateral-force range.
Each rollout records per-frame chassis position, yaw, speed, yaw rate, wheel angular velocities, and steer angles at 60\,Hz.

\paragraph{Loss function.}
Let $x^{\mathrm{ref}}_{k,t}$ denote the teacher's value for channel $k$ at frame $t$, and $x^{\mathrm{stu}}_{k,t}$ the corresponding student value under a candidate parameter vector.
Each rollout has $T$ frames recorded at 60\,Hz.
The per-channel mean squared error is $\mathrm{MSE}_k = \tfrac{1}{T}\sum_{t=1}^{T}(x^{\mathrm{stu}}_{k,t} - x^{\mathrm{ref}}_{k,t})^2$.
The full objective adds two terminal-frame penalties that emphasize end-of-maneuver accuracy:
\begin{equation}
\label{eq:app_sysid_loss}
\mathcal{L} = \sum_{k \in \mathcal{K}} w_k \cdot \mathrm{MSE}_k
  + w_{\mathrm{pos,final}} \cdot \|\mathbf{p}^{\mathrm{stu}}_T - \mathbf{p}^{\mathrm{ref}}_T\|^2
  + w_{\mathrm{spd,final}} \cdot (s^{\mathrm{stu}}_T - s^{\mathrm{ref}}_T)^2,
\end{equation}
where $\mathcal{K} = \{\text{position}_{xy}, \text{yaw}, \text{speed}, \text{yaw rate}, \text{wheel speed}, \text{steer angle}\}$ with weights $w_k = (1.0, 0.4, 0.35, 0.25, 0.05, 0.05)$ respectively.
The vectors $\mathbf{p}^{\mathrm{stu}}_T, \mathbf{p}^{\mathrm{ref}}_T \in \mathbb{R}^2$ are the student and teacher XY positions at the final frame, and $s^{\mathrm{stu}}_T, s^{\mathrm{ref}}_T$ are the corresponding speeds.
The terminal weights are $w_{\mathrm{pos,final}} = 1.5$ and $w_{\mathrm{spd,final}} = 0.75$.

\paragraph{Multi-stage CEM search.}
The CEM search is divided into five stages to manage the 20-dimensional space:
\begin{enumerate}[leftmargin=*, itemsep=1pt]
  \item \emph{Longitudinal}: tunes drive/brake torques and wheel parameters on the straight-line rollout.
  \item \emph{Steering}: tunes steering gains, suspension, and damping on the turning rollouts.
  \item \emph{Surface}: tunes per-surface friction scales on the transition rollout.
  \item \emph{Refinement}: tunes all 20 parameters jointly on all rollouts with a narrowed search window (18\% of the full range).
  \item \emph{Brake preservation}: re-tunes longitudinal parameters with a tight window (10\%) to prevent regression.
\end{enumerate}
Each stage inherits the best configuration from the previous stage.
The CEM uses a population of 24 candidates, an elite fraction of 25\%, an initial standard deviation of 25\% of each parameter's range, and a minimum standard deviation of 5\% to prevent premature convergence.
A total of 320 trials are allocated across the five stages with weights (30\%, 20\%, 15\%, 20\%, 15\%).

\paragraph{Vehicle-intrinsic vs.\ environment parameters.}
The 20 sysid parameters are \emph{vehicle-intrinsic}: they characterize how the articulation converts actuator commands into motion (drive torque, steering response, suspension compliance, etc.).
They do not change when the road surface changes.
Environment-level friction variation is handled orthogonally by the PhysX contact solver: each environment's wheel colliders are assigned a friction coefficient $\mu$ via their physics material, and the contact solver computes friction forces at the wheel--ground contact points accordingly (using combine mode \texttt{min}).
This separation is deliberate: sysid produces a realistic vehicle model once, and the physics engine then automatically adapts that vehicle's behavior to any friction environment.

\chattodo{\textbf{ACTION ITEM --- FIGURE:} Consider a trajectory overlay figure showing teacher vs.\ student vehicle on one or two maneuvers.}


% ============================================================
\section{Simulator Limitations and Sharp Edges}
\label{app:sharp_edges}
% ============================================================

Transparency about simulator assumptions is essential for interpreting the experimental results.
Below we catalogue the most consequential simplifications, known failure modes, and practical workarounds in SceneFactory, organized by subsystem.

\paragraph{Z-flattening.}
The Waymo Open Motion Dataset provides full 3D polyline coordinates, including road elevation changes across bridges, ramps, overpasses, and multi-level interchanges.
SceneFactory collapses all road geometry to a constant ground height ($z{=}0$), because vehicles drive on a single flat PhysX ground plane per world.
This design decision has two consequences.
First, scenes with elevated structures (bridges, highway overpasses) produce visually overlapping road segments at ground level, which can confuse the $k$-nearest-road-point observation by presenting the agent with road points from an overhead or underpass road that it cannot physically reach.
Second, ramp grades and banked curves are lost, removing a class of driving challenges (hill starts, crest reactions) that affect real-world vehicle dynamics.
We mitigate the first issue through scene curation and by filtering road points whose original Waymo $z$ coordinate deviates significantly from the ground plane, but the fundamental limitation remains: SceneFactory is a 2.5D simulator operating on 3D data.

\paragraph{Scene curation.}
Not all Waymo scenarios are suitable for flat-ground training.
We developed an interactive \emph{Config Wizard} tool (\texttt{scene\_factory\_config\_wizard.py}) that presents each candidate scene as a bird's-eye plot showing lane-center polylines (Waymo types 1--2, blue), road-edge polylines (types 15--16, red), and valid start--goal pairs connected by origin--destination lines.
A human reviewer accepts or rejects each scene via keyboard input (y/n/b/q), with undo support.
The tool supports both a lightweight Matplotlib mode for rapid triage and an Isaac Sim preview mode that renders the full USD road geometry with 3D start/goal markers.

From a pool of ${\sim}100$ converted Waymo scenarios, we curated 64 training scenes and 64 held-out evaluation scenes, rejecting 23 training candidates and 49 evaluation candidates.
Common rejection reasons include: (i)~multi-level interchanges where z-flattening produces overlapping geometry that would poison the road-point observation, (ii)~scenes where origin--destination pairs straddle a road-edge polyline, making the lane-forbidden penalty unavoidable, and (iii)~scenes with fewer than two valid controllable agents after spawn filtering.
The wizard produces a complete training configuration (scene list, environment count, reward weights, and curriculum stage) from scratch, so a researcher can go from raw Waymo data to a reproducible training run with minimal manual editing.

\paragraph{Waymo road-edge labeling.}
SceneFactory treats Waymo polyline types 15 and 16 as forbidden road edges for both the lane-forbidden termination event and the road-edge TTC penalty.
In practice, these labels are not uniformly reliable across the dataset.
As noted by Kazemkhani et al.~\cite{kazemkhani_gpudrive_2025}, the Waymo dataset labels parking-lot entrances as road edges.
Some tracked vehicles in these scenarios have their origin on the main road and their destination inside the parking area (or vice versa), so the recorded trajectory necessarily crosses a road-edge polyline.
When such a scenario is used for training, the agent faces an impossible conflict: reaching the goal requires crossing a boundary that triggers the lane-forbidden penalty ($-20.0$) and immediate termination.
These cases inject contradictory supervision into the training signal, degrading both goal-reach rate and lane-safety metrics.
Other scenarios lack road-edge annotations entirely where a physical curb or barrier clearly exists, creating blind spots in the lane-safety reward.
Scene curation removes the worst cases, but residual label noise remains and likely contributes to the modest lane-forbidden violation rates observed in evaluation.

\paragraph{Moving-only agent assumption.}
The spawn filter retains only Waymo agents whose start and goal positions are separated by more than 3.0\,m (the goal-reached threshold).
Agents below this threshold are treated as ``already at goal'' and excluded entirely.
Because the Waymo dataset records start and end positions from the logged trajectory, vehicles that remain parked, idle, or nearly stationary throughout the 9.1\,s recording window have near-zero start--goal displacement and are filtered out.
As a result, \emph{parked vehicles do not exist in the simulation}: a scene in which a moving vehicle must navigate past a row of curbside-parked cars will contain only the moving vehicle.
This simplification removes an entire class of realistic interactions (door-zone avoidance, double-parked obstacle negotiation, parallel parking maneuvers) and inflates the effective lane width available to the agent.
A future extension could retain filtered agents as static rigid-body obstacles rather than discarding them.

\paragraph{PointInstancer instability under Fabric.}
Road segments are rendered as USD \texttt{PointInstancer} prims for compactness (one instancer per road type per world, with per-instance position, orientation, and scale arrays).
We discovered that PointInstancer prims interact poorly with Isaac Lab's Fabric state management layer: visibility updates and prototype-table initialization are unreliable, producing intermittent rendering artifacts (invisible or misplaced road segments) that do not affect physics but corrupt visual debugging and video captures.
An explicit-prim fallback mode is provided (one \texttt{UsdGeom.Cube} per road segment) for diagnostic and stage-export workflows.
Production training disables Fabric entirely (\texttt{use\_fabric: false}) to avoid related issues with road metadata access, at a modest throughput cost.

\paragraph{Collision-sensor warmup.}
When vehicles are spawned or teleported to their start poses, PhysX requires several substeps to settle the suspension and establish stable wheel--ground contact.
During this transient, contact sensors report large spurious forces as the chassis drops onto the ground plane.
Without mitigation, these forces trigger false collision terminations at the start of every episode.
We suppress collision detection for the first 24 environment steps (${\approx}0.8$\,s) after each spawn, which is sufficient for the suspension to settle.
This warmup window introduces a brief period at episode start during which actual inter-vehicle collisions are unpenalized, though in practice vehicles are spawned with sufficient separation that early collisions are rare.

\paragraph{Teleport-only reset.}
Isaac Lab's default environment reset tears down and reconstructs PhysX actors, which is prohibitively expensive for multi-agent scenes with per-environment road geometry.
SceneFactory uses a \emph{teleport-only} reset strategy: after a single full initialization reset, subsequent episode resets reposition vehicles by writing directly to root-state tensors (position, orientation, linear and angular velocity) without rebuilding the physics scene.
This is ${\sim}10{\times}$ faster but has a subtle limitation: accumulated PhysX internal state (joint impulse caches, contact-pair broadphase data) is not fully cleared.
In rare cases this causes a vehicle to retain a residual velocity or contact state from the previous episode for one or two substeps after teleport.
The collision warmup window (above) masks most such artifacts.

\paragraph{Uniform friction surface.}
Per-world friction variation is implemented by writing to each vehicle's wheel collider physics material at environment construction time.
Because the global ground plane has friction $\mu{=}1.0$ and PhysX uses combine mode \texttt{min}, the effective contact friction is $\min(\mu_{\mathrm{wheel}}, 1.0) = \mu_{\mathrm{wheel}}$, giving per-environment control.
However, friction remains spatially uniform within each world: every wheel contact with the ground plane uses the same coefficient.
Real roads exhibit local friction variation (wet patches, oil spills, transitions between asphalt and concrete), and scenarios such as a slippery lane adjacent to a dry lane cannot be represented.
Modeling local friction would require either a segmented ground mesh with per-face materials or runtime material-property writes keyed to vehicle position, neither of which is straightforward under PhysX~5's current contact pipeline.

\paragraph{PhysX substep granularity.}
Physics runs at 120\,Hz with a decimation factor of 4, yielding 30\,Hz control.
The 120\,Hz rate is a compromise between simulation fidelity and throughput: higher rates (240\,Hz or 480\,Hz) would improve suspension stability and contact resolution but proportionally reduce training speed.
Our sysid calibration (\S\ref{sec:vehicle_sysid}) was performed at this rate, so the fitted parameters absorb any 120\,Hz-specific artifacts; re-calibration would be needed if the substep rate were changed.

\paragraph{Clone-in-Fabric disabled.}
Isaac Lab's \texttt{clone\_in\_fabric} optimization, which clones environment state entirely in the Fabric layer for maximum GPU throughput, is disabled in SceneFactory because it is incompatible with the post-clone road injection workflow described in \S\ref{sec:data_processing}.
Fabric assumes that all environments are structurally identical after cloning; SceneFactory's per-environment road prims violate this assumption, causing metadata reads to return stale or incorrect values.
Disabling Fabric increases CPU--GPU synchronization overhead but is necessary for correctness.

\paragraph{Constant-velocity TTC.}
Both the inter-vehicle and road-edge TTC computations assume constant velocity over the prediction horizon.
This systematically overestimates collision risk during deceleration (the vehicle is predicted to continue at its current speed toward the obstacle) and underestimates risk during acceleration.
A quadratic (constant-acceleration) TTC formulation would be more accurate, particularly in approach zones where vehicles are actively braking, but would increase the per-step computational cost of the $K_v{\times}9$ pairwise circle tests.


% ============================================================
\section{Throughput Analysis Details}
\label{app:throughput}
% ============================================================

\paragraph{Timing breakdown.}
Table~\ref{tab:app_timing_breakdown} reports the per-step timing breakdown for both the Baseline and SceneFactory pipelines at 64 worlds (1{,}024 agent slots, all active), from the Experiment~A scaling benchmark on scene \texttt{000077}.

On the Baseline side, the PhysX rigid-body solver accounts for roughly half of wall-clock time (2{,}934\,ms, 51.6\%); this cost is irreducible and shared by both pipelines.
The observation builder consumes another 1{,}456\,ms (25.6\%), because it iterates over each of the 1{,}024 active agents in a Python loop, issuing per-prim USD state queries and computing ego-frame transforms with scalar arithmetic.
Geometric lane features, TTC computation, and action application add a further 22\% through similar per-agent loops.
The total per-step cost is \textbf{5{,}683\,ms}.

On the SceneFactory side, the same 64-world configuration completes an environment step in \textbf{135\,ms}---a \textbf{42$\times$} reduction.
Physics remains the largest single cost (53.8\,ms), but the GPU-batched solver is 55$\times$ faster than the Baseline's single-scene PhysX dispatch.
Observations, rewards, dones, and actions collectively take only 78\,ms because they execute as batched \texttt{torch} tensor operations with no per-agent Python loop.

\begin{table}[t]
\centering
\caption{Per-step timing breakdown at 64 worlds (1{,}024 agent slots, all active) from the Experiment~A benchmark on scene \texttt{000077}.
Baseline sub-phases are grouped into the matching SceneFactory category.
SceneFactory sub-phases may overlap on the GPU, so they sum to more than the measured wall-clock total; the \textbf{Total} row reports the true measured time.}
\label{tab:app_timing_breakdown}
\small
\begin{tabular}{lrrr}
\toprule
Phase & Baseline (ms) & SceneFactory (ms) & Speedup \\
\midrule
Physics solver         & 2{,}934 &  53.8 & 55$\times$ \\
Observation builder    & 1{,}772 &  23.3 & 76$\times$ \\
Reward \& termination  &     629 &  29.9 & 21$\times$ \\
Action application     &     299 &  24.9 & 12$\times$ \\
Reset / overhead       &      49 &  32.4 &  1.5$\times$ \\
\midrule
\textbf{Total (wall-clock)} & \textbf{5{,}683} & \textbf{135} & \textbf{42$\times$} \\
\bottomrule
\end{tabular}
\\[4pt]
{\footnotesize
\emph{Baseline grouping:}
Observation builder = obs.\ construction (1{,}456\,ms) + geometric lane features (316\,ms);
Reward \& termination = road-edge TTC (236) + vehicle TTC (191) + TTC state queries (136) + contact scan (66).
\emph{SceneFactory grouping:}
Reward \& termination = reward computation (18.4) + done/termination (11.5);
Physics solver = write (33.0) + simulate (15.9) + read-back (4.9).}
\end{table}

\paragraph{Why the Baseline is slow.}
SceneFactory eliminates Baseline bottlenecks through three architectural changes:
\begin{enumerate}[leftmargin=*, itemsep=2pt]
  \item \textbf{Batched joint-state tensors.}
        The custom URDF articulated vehicle (\S\ref{sec:vehicle_sysid}) exposes positions, velocities, and orientations as GPU-resident tensors of shape $(\texttt{num\_envs}, \texttt{num\_agents}, d)$ via Isaac Lab's \texttt{Articulation} API.
        Observation construction reduces to a handful of batched \texttt{torch} operations with no Python per-agent loop.
  \item \textbf{Cloned PhysX scenes.}
        Isaac Lab's environment cloning creates $N$ independent PhysX solver islands, enabling the GPU broad-phase and narrow-phase to dispatch them in parallel.
        The Baseline places all worlds in a single PhysX scene, limiting solver parallelism.
  \item \textbf{GPU-resident training loop.}
        RSL-RL reads observations, rewards, and dones as CUDA tensors directly from the environment; no CPU round-trip occurs during training.
        The Baseline converts observations to NumPy for Stable-Baselines3's rollout buffer, incurring $O(\text{agents} \times \text{obs\_dim})$ bytes of device-to-host transfer every step.
\end{enumerate}


% ============================================================
\section{PPO Hyperparameters}
\label{app:ppo_hparams}
% ============================================================

Table~\ref{tab:app_ppo_hparams} lists the full PPO configuration used in all experiments.
The learning rate follows an adaptive schedule that halves or doubles the rate to maintain the per-update KL divergence near a target of 0.01.

\begin{table}[h]
\centering
\caption{PPO hyperparameters.}
\label{tab:app_ppo_hparams}
\small
\begin{tabular}{ll}
\toprule
Parameter & Value \\
\midrule
Rollout horizon           & 128 steps \\
Mini-batches              & 16 \\
Learning epochs           & 5 \\
Learning rate (initial)   & $2.09 \times 10^{-4}$ \\
LR schedule               & adaptive (KL target 0.01) \\
Discount $\gamma$         & 0.99 \\
GAE $\lambda$             & 0.98 \\
Clip $\epsilon$           & 0.15 \\
Entropy coefficient       & $3 \times 10^{-5}$ \\
Value loss coefficient    & 1.0 \\
Max gradient norm         & 1.0 \\
\bottomrule
\end{tabular}
\end{table}


\end{document}
